\documentclass[a4paper,11pt]{article}

\usepackage[utf8]{inputenc}
\usepackage[T1]{fontenc}
\usepackage{comment}
\usepackage[margin=2.2cm, top=2.8cm, bottom=2.8cm]{geometry}
\usepackage{xcolor}
\usepackage{geometry}
\usepackage{caption}
\usepackage{float}
\usepackage{graphicx}
\usepackage{float}
\usepackage{titlesec}
\usepackage{enumitem}
\usepackage{fancyhdr}
\usepackage{booktabs}
\usepackage{tabularx}
\usepackage{array}
\usepackage[hidelinks]{hyperref}
\usepackage{tcolorbox}
\usepackage{parskip}
\usepackage{mdframed}
\usepackage{multicol}
\usepackage{amsmath}
\usepackage{amssymb}
\usepackage{bm}
\usepackage{tikz}
\usetikzlibrary{positioning,shapes.geometric,arrows.meta,fit,backgrounds,calc,bending}

\tcbuselibrary{skins, breakable}

% --- Couleurs -------------------------------------------------------------------
\definecolor{navy}{HTML}{0D2A4A}
\definecolor{cblue}{HTML}{1A56A0}
\definecolor{lightblue}{HTML}{D6E4F5}
\definecolor{cred}{HTML}{C0392B}
\definecolor{redbg}{HTML}{FAE5E3}
\definecolor{cgray}{HTML}{555555}
\definecolor{lightgray}{HTML}{F5F5F5}
\definecolor{gold}{HTML}{9A7000}
\definecolor{goldbg}{HTML}{FFFBEA}
\definecolor{green}{HTML}{1A6B35}
\definecolor{greenbg}{HTML}{E8F5EE}
\definecolor{codeinline}{HTML}{2C3E80}
\definecolor{riskgreen}{HTML}{1A6B35}
\definecolor{purple}{HTML}{6A0DAD}
\definecolor{purplebg}{HTML}{F3EAFE}

% --- Titres ---------------------------------------------------------------------
\titleformat{\section}{\Large\bfseries\color{navy}}{}{0em}{}[\color{cblue}\titlerule]
\titleformat{\subsection}{\large\bfseries\color{cblue}}{}{0em}{}
\titleformat{\subsubsection}{\normalsize\bfseries\color{cgray}}{}{0em}{}
\titlespacing{\section}{0pt}{1.8em}{0.6em}
\titlespacing{\subsection}{0pt}{1.2em}{0.4em}
\titlespacing{\subsubsection}{0pt}{0.8em}{0.2em}

% --- Inline code ----------------------------------------------------------------
\newcommand{\cs}[1]{\texttt{\color{codeinline}#1}}

% --- Boites ---------------------------------------------------------------------
\tcbset{
  constraint/.style={
    breakable, colback=redbg, colframe=cred,
    colbacktitle=cred, coltitle=white,
    fonttitle=\bfseries\small,
    title={Contraintes C\# -- NE PAS IGNORER},
    left=8pt, right=8pt, top=6pt, bottom=6pt,
    boxrule=0.8pt, arc=3pt,
  },
  note/.style={
    breakable, colback=goldbg, colframe=gold,
    colbacktitle=gold, coltitle=white,
    fonttitle=\bfseries\small,
    title={Pourquoi c'est important},
    left=8pt, right=8pt, top=6pt, bottom=6pt,
    boxrule=0.8pt, arc=3pt,
  },
  deliverable/.style={
    breakable, colback=greenbg, colframe=green,
    colbacktitle=green, coltitle=white,
    fonttitle=\bfseries\small,
    title={Crit\`ere de validation},
    left=8pt, right=8pt, top=6pt, bottom=6pt,
    boxrule=0.8pt, arc=3pt,
  },
  math/.style={
    breakable, colback=purplebg, colframe=purple,
    colbacktitle=purple, coltitle=white,
    fonttitle=\bfseries\small,
    title={Mod\'elisation math\'ematique},
    left=8pt, right=8pt, top=6pt, bottom=6pt,
    boxrule=0.8pt, arc=3pt,
  },
  todo/.style={
    breakable, colback=lightgray, colframe=cgray,
    colbacktitle=cgray, coltitle=white,
    fonttitle=\bfseries\small,
    title={TO BE CONTINUED -- impl\'ementation ult\'erieure},
    left=8pt, right=8pt, top=6pt, bottom=6pt,
    boxrule=0.8pt, arc=3pt,
  },
}

% --- Header / Footer ------------------------------------------------------------
\pagestyle{fancy}
\fancyhf{}
\fancyhead[L]{\small\color{navy}\textbf{Trading Dashboard -- C\# .NET 8}}
\fancyhead[R]{\small\color{cgray}\textit{Sp\'ecifications Techniques v2.0}}
\fancyfoot[C]{\small\color{cgray}\thepage}
\fancyfoot[R]{\small\color{cgray}\textit{open-source -- github.com}}
\renewcommand{\headrulewidth}{0.4pt}
\renewcommand{\footrulewidth}{0.2pt}

% Macros maths
\newcommand{\E}{\mathbb{E}}
\newcommand{\Q}{\mathbb{Q}}
\newcommand{\R}{\mathbb{R}}
\newcommand{\dd}{\mathrm{d}}

% ============================================================
\begin{document}

% --- PAGE DE GARDE --------------------------------------------------------------
\begin{titlepage}
\centering
\vspace*{2cm}
{\huge\bfseries\color{navy} Trading Dashboard\par}

\vspace{1.2cm}
\noindent\rule{\textwidth}{1.5pt}
\vspace{0.5cm}
{\LARGE\bfseries\color{navy} Sp\'ecifications Techniques\par}
\vspace{0.3cm}
{\large\color{cgray} Guide d'impl\'ementation -- C\# .NET 8\par}
\vspace{0.5cm}
\noindent\rule{\textwidth}{1.5pt}



\vspace{3cm}
\begin{tabular}{ll}
\toprule
\textbf{Stack}    & C\# 12, .NET 8 LTS, Entity Framework Core 8 \\
\textbf{Version}  & 2.0 -- 2026 \\
\textbf{Modules}  & 8 projets \\
\textbf{Tests}    & xUnit, Moq, FluentAssertions \\
\midrule
\textbf{Nouveaut\'es v2.0} & Architecture en couches path/payoff/MC \\
& Produits : barrières, digitales, worst-of, variance swaps ...\\
& Pricer: multi-devises/multi-actifs/multi-flux \\
& Model: Heston, LV, LSV \\
\bottomrule
\end{tabular}
\end{titlepage}

% --- TABLE DES MATIERES ---------------------------------------------------------
\tableofcontents
\newpage

% ============================================================
\section{Vue d'ensemble du projet}
% ============================================================

\subsection{Objectif et p\'erim\`etre produits}

\textbf{Trading Dashboard} est un syst\`eme de pricing et de gestion de risque en temps
quasi-r\'eel, impl\'ement\'e en C\# .NET 8. Il couvre l'ensemble de la cha\^ine de valeur
d'un portefeuille multi-actifs, multi-devises : valorisation, Greeks, PnL, limites
de risque, stress testing et rapports.

La version 2.0 \'etend consid\'erablement le p\'erim\`etre produits par rapport \`a la v1.2
(vanilles et exotiques mono-actif). Toutes les nouvelles familles sont priceable sous
Black-Scholes Monte Carlo dans un premier temps~; les mod\`eles \`a volatilit\'e stochastique
(Heston, Local Vol, LSV) feront l'objet d'it\'erations ult\'erieures.

\begin{center}
\begin{tabularx}{\textwidth}{l X l}
\toprule
\textbf{Famille} & \textbf{Produits} & \textbf{Moteur v2.0} \\
\midrule
Vanilles mono-actif & Call, Put, Future, Forward & BS analytique \\
Exotiques mono-actif & AsianCall/Put, LookbackCall/Put (fixe et flottant) & MC mono-actif \\
Barrières mono-actif & Knock-In/Out (Up/Down), Double Barrier & MC mono-actif + BGK \\
Digitales mono-actif & CashOrNothing, AssetOrNothing & MC mono-actif \\
FX Options & FxCall, FxPut, FxAsian, FxBarrier & MC multi-devises \\
Multi-actifs & BasketCall/Put, WorstOf, BestOf, Spread, Outperformance & MC multi-actifs \\
Quanto & Quanto (actif \'etranger, paiement domestique) & MC multi-devises \\
Variance / Volatilit\'e & VarianceSwap, VolatilitySwap & MC mono-actif \\
Cliquets & Cliquet, Ratchet & MC mono-actif \\
Structured (placeholder) & AutocallAthena, AutocallPhoenix, ReverseConvertible & \textit{To be continued} \\
\bottomrule
\end{tabularx}
\end{center}

\begin{tcolorbox}[note]
\textbf{Principe fondamental du choix de mod\`ele.} La question n'est pas
``formule ferm\'ee ou Monte Carlo ?'' mais ``le mod\`ele capture-t-il les risques
d\'eterminants du prix de march\'e ?'' Sous Black-Scholes, les barrières ont une
formule ferm\'ee (Rubinstein-Reiner) mais elle ignore le skew de volatilit\'e, qui
est pr\'ecis\'ement maximal \`a hauteur de la barri\`ere. En prod, les barrières sont
pric\'ees sous Local Volatility ou LSV calibr\'e \`a la surface compl\`ete. La version BS
de ce projet constitue un pricer de r\'ef\'erence, non un pricer de march\'e.
\end{tcolorbox}

\subsection{Architecture en couches -- principe directeur}

La version 2.0 repose sur \textbf{trois axes d'extension ind\'ependants} :

\begin{enumerate}[leftmargin=1.5em]
  \item \textbf{Le mod\`ele} (BS, Heston, LV, LSV) -- d\'etermine comment on g\'en\`ere les
        trajectoires. Extension point : \cs{IPathGenerator}.
  \item \textbf{Le produit} (call, barri\`ere, basket, autocall) -- d\'etermine comment
        on calcule le payoff. Extension point : \cs{IPayoff}.
  \item \textbf{La technique MC} (brute force, pont brownien, variable de contr\^ole) --
        d\'etermine comment on estime l'esp\'erance. Extension point : \cs{IVarianceReduction}.
\end{enumerate}

Ces trois axes ne se contaminent pas. Ajouter Heston ne modifie ni les payoffs ni
le moteur MC. Ajouter un autocall ne modifie ni le simulateur ni le moteur.

\textbf{Couches du module Pricing :}

\begin{center}
\begin{tabularx}{\textwidth}{l l X}
\toprule
\textbf{Couche} & \textbf{Abstraction principale} & \textbf{R\^ole} \\
\midrule
1 -- March\'e & \cs{MarketEnvironment} & Toutes les donn\'ees de march\'e : spots, taux, FX, matrice de corr\'elation \\
2 -- Paths & \cs{IPathGenerator} / \cs{SimulatedPaths} & G\'en\'eration de trajectoires, ind\'ependante du produit \\
3 -- Payoffs & \cs{IPayoff} & Calcul du flux non actualis\'e depuis les trajectoires \\
4 -- R\'eduction de variance & \cs{IVarianceReduction} & Pont brownien, variables de contr\^ole, antith\'etique \\
5 -- Moteur MC & \cs{IMonteCarloEngine} & Estimation de l'esp\'erance, actualisation, IC \`a 95\% \\
6 -- Portefeuille & \cs{PortfolioPricer} & Simulation conjointe de toutes les positions \\
\bottomrule
\end{tabularx}
\end{center}

\subsection{Architecture de la solution (inchang\'ee)}

La solution \texttt{TradingDashboard.sln} contient 8 projets. Graphe de d\'ependances
identique \`a la v1.2 : les d\'ependances sont unidirectionnelles, aucune r\'ef\'erence
circulaire.

\begin{center}
\begin{tabularx}{\textwidth}{>{\ttfamily}l l X}
\toprule
\textbf{Projet} & \textbf{Type} & \textbf{Responsabilit\'e} \\
\midrule
.Core       & Class Library  & Entit\'es, interfaces, enums. Aucune d\'ependance externe. \\
.Pricing    & Class Library  & BS analytique ; Monte Carlo en couches (IPathGenerator,
                               IPayoff, IVarianceReduction, IMonteCarloEngine) ;
                               PortfolioPricer ; Greeks ; Factory/Strategy. \\
.Data       & Class Library  & Entity Framework Core, migrations, Repository. \\
.MarketData & Class Library  & Alpha Vantage, cache, retry Polly. \\
.Risk       & Class Library  & PnL Engine, LimitChecker, StressTest, VaR/CVaR
                               avec simulation conjointe. \\
.Reports    & Class Library  & Export Excel (EPPlus) et PDF (QuestPDF). \\
.Dashboard  & Console App    & Point d'entr\'ee, DI, affichage temps r\'eel. \\
.Tests      & xUnit Project  & Tests unitaires et d'int\'egration. \\
\bottomrule
\end{tabularx}
\end{center}

\subsection{S\'equence de d\'eveloppement v2.0}

\begin{center}
\begin{tabularx}{\textwidth}{l l X}
\toprule
\textbf{Semaine} & \textbf{Module} & \textbf{Objectif} \\
\midrule
1   & .Core v2.0       & ProductType \'etendu, IPathGenerator, SimulatedPaths, IPayoff v2, IVarianceReduction \\
2   & .Pricing -- Paths & GbmPathGenerator, MultiAssetGbmPathGenerator (Cholesky), MultiCurrencyGbmPathGenerator \\
3   & .Pricing -- Payoffs mono-actif & Barrières (BGK), Digitales, Lookback flottant, VarianceSwap \\
4   & .Pricing -- Payoffs multi-actifs & BasketPayoff, WorstOfPayoff, BestOfPayoff, SpreadPayoff \\
5   & .Pricing -- Moteur MC & IMonteCarloEngine, variance reduction, PortfolioPricer \\
6   & .Risk v2.0       & PortfolioSimulator conjoint, PnLEngine \\
7   & .Risk -- Limites \& Stress & LimitChecker, StressTestEngine \\
8   & .Reports + .Dashboard & Export, affichage, DI \\
\bottomrule
\end{tabularx}
\end{center}

\newpage
% ============================================================
\section{Module Core v2.0 -- \texttt{TradingDashboard.Core}}
% ============================================================

\subsection{Enums}

\subsubsection{\texttt{ProductType.cs} -- Version 2.0}

L'enum couvre d\'esormais l'ensemble des familles de produits d\'ecrites en section 1.

\textbf{Valeurs attendues (avec attribut \cs{[Description]} sur chacune) :}
\begin{center}
\begin{tabularx}{1.1\textwidth}{>{\ttfamily}l l l}
\toprule
\textbf{Valeur} & \textbf{Famille} & \textbf{Description m\'etier} \\
\midrule
Call/Put & Vanilla mono-actif & European Call / Put Option \\
Future/Forward & Vanilla mono-actif & Futures / Forward Contract \\
AsianCall/Put & Exotique mono-actif & Asian Arithmetic Average Call / Put \\
LookbackCall/Put & Exotique mono-actif & Lookback Fixed Strike Call / Put \\
FloatingLookbackCall/Put & Exotique mono-actif & Lookback Floating Strike Call / Put \\
UpAndOutCall/Put & Barri\`ere knock-out & Up-and-Out / Down-and-Out Call / Put \\
DownAndOutCall/Put & Barri\`ere knock-out & Down-and-Out Call / Put \\
UpAndInCall/Put & Barri\`ere knock-in & Up-and-In Call / Put \\
DownAndInCall/Put & Barri\`ere knock-in & Down-and-In Call / Put \\
DoubleBarrierCall/Put & Double barri\`ere & Double Barrier Call / Put \\
CashOrNothingCall/Put & Digitale & Cash-or-Nothing Call / Put \\
AssetOrNothingCall/Put & Digitale & Asset-or-Nothing Call / Put \\
FxCall/Put & FX Options & FX European Call / Put (Garman-Kohlhagen) \\
FxAsianCall/Put & FX Exotique & FX Asian Arithmetic Average Call / Put \\
FxBarrierCall/Put & FX Barri\`ere & FX Knock-Out Call / Put \\
BasketCall/Put & Multi-actifs & Basket Call / Put (moyenne pond\'er\'ee) \\
WorstOfCall/Put & Multi-actifs & Worst-of Call / Put (minimum des rendements) \\
BestOfCall/Put & Multi-actifs & Best-of Call / Put (maximum des rendements) \\
SpreadOption & Multi-actifs & Spread Option ($S^1 - S^2 - K$) \\
Outperformance & Multi-actifs & Outperformance ($S^1 - S^2$, $K=0$) \\
Quanto & Multi-devises & Quanto (actif \'etranger, paiement domestique) \\
VarianceSwap & Variance / Vol & Variance Swap \\
VolatilitySwap & Variance / Vol & Volatility Swap \\
Cliquet & Cliquet & Cliquet (resetted-strike sum of returns) \\
Ratchet & Cliquet & Ratchet (accumulated capped returns) \\
AutocallAthena & Structured (placeholder) & Autocall Athena \\
AutocallPhoenix & Structured (placeholder) & Autocall Phoenix \\
ReverseConvertible & Structured (placeholder) & Reverse Convertible \\
\bottomrule
\end{tabularx}
\end{center}

\textbf{M\'ethodes d'extension dans \cs{ProductTypeExtensions} :}

\begin{itemize}[leftmargin=1.5em]
  \item \cs{IsVanilla()} : Call, Put uniquement.
  \item \cs{IsExotic()} : toutes les familles mono-actif path-dependent (Asian, Lookback, Barrier, Digital, Variance, Cliquet).
  \item \cs{IsOption()} : tout produit avec optionnalit\'e (hors Future, Forward).
  \item \cs{IsFx()} : FxCall, FxPut, FxAsianCall, FxAsianPut, FxBarrierCall, FxBarrierPut, Quanto.
  \item \cs{IsMultiAsset()} : Basket, WorstOf, BestOf, Spread, Outperformance.
  \item \cs{IsBarrier()} : les 10 types barri\`eres (knock-in, knock-out, double).
  \item \cs{IsDigital()} : CashOrNothing et AssetOrNothing.
  \item \cs{IsStructured()} : AutocallAthena, AutocallPhoenix, ReverseConvertible.
  \item \cs{RequiresMultipleUnderlyings()} : \cs{IsMultiAsset() || IsStructured()}.
\end{itemize}

\begin{tcolorbox}[constraint]
\cs{IsVanilla()}, \cs{IsExotic()}, \cs{IsFx()}, \cs{IsMultiAsset()} sont mutuellement
exclusifs. Un switch exhaustif dans \cs{PricerFactory} doit couvrir tous les cas
sans \cs{default} silencieux -- utiliser \cs{_ => throw new UnsupportedProductTypeException(type)}.
\end{tcolorbox}

\subsubsection{\texttt{AssetClass.cs}, \texttt{Currency.cs}, \texttt{AlertLevel.cs}}

Inchang\'es par rapport \`a la v1.2. \cs{AssetClass} comporte \cs{Equity}, \cs{FX}, \cs{Rate}.

\subsection{Entit\'es}

\subsubsection{\texttt{MarketEnvironment.cs} -- Remplacement de MarketSnapshot}

\cs{MarketEnvironment} est l'objet central de march\'e pass\'e \`a tous les composants
du moteur de pricing v2.0.

\textbf{Propri\'et\'es attendues :}
\begin{itemize}[leftmargin=1.5em]
  \item \cs{double DomesticRate} -- taux sans risque domestique $r^0$.
  \item \cs{IReadOnlyDictionary<Currency, double> ForeignRates} -- $r^i$ par devise.
  \item \cs{IReadOnlyDictionary<string, MarketSnapshot> Snapshots} -- snapshot par ticker.
  \item \cs{IReadOnlyDictionary<(Currency, Currency), FxRate> FxRates} -- taux de change.
  \item \cs{GlobalCorrelationMatrix Correlation} -- matrice de corr\'elation globale.
  \item \cs{DateTimeOffset AsOf} -- horodatage de l'environnement.
\end{itemize}

\textbf{Factory m\'ethods (r\'etrocompatibilit\'e v1.2) :}
\begin{itemize}[leftmargin=1.5em]
  \item \cs{static MarketEnvironment FromSingleAsset(MarketSnapshot snap)} -- construit
        un environnement mono-actif domestique, permettant aux anciens pricers
        de continuer \`a fonctionner sans modification.
  \item \cs{static MarketEnvironment FromSnapshots(IEnumerable<MarketSnapshot> snaps,}\\
        \cs{double domesticRate, GlobalCorrelationMatrix corr)} -- construction multi-actifs.
\end{itemize}

\subsubsection{\texttt{GlobalCorrelationMatrix.cs}}

\textbf{Structure :} matrice carr\'ee \cs{double[,]} de taille $N \times N$, o\`u $N = \sum n_i + \bar{n}$.
L'ordonnancement des facteurs suit exactement la convention math\'ematique d\'efinie
en section~3.1.

\textbf{Propri\'et\'es :}
\begin{itemize}[leftmargin=1.5em]
  \item \cs{double[,] Rho} -- la matrice de corr\'elation.
  \item \cs{AssetIndex[] AssetMapping} -- mapping index $\to$ (ticker, AssetClass, Currency).
  \item \cs{double[,] CholeskyFactor} -- facteur $L$ tel que $LL^\top = \rho$, calcul\'e \`a la construction.
\end{itemize}

\textbf{M\'ethodes :}
\begin{itemize}[leftmargin=1.5em]
  \item \cs{void Validate()} -- v\'erifie sym\'etrie, 1 sur la diagonale, d\'efinie positive.
  \item \cs{int IndexOf(string ticker, AssetClass cls)} -- r\'ecup\`ere l'index dans le vecteur global.
  \item \cs{GlobalCorrelationMatrix Sub(int[] indices)} -- sous-matrice pour un sous-ensemble d'actifs.
\end{itemize}

\begin{tcolorbox}[note]
La d\'ecomposition de Cholesky est calcul\'ee une seule fois \`a la construction. Elle est
r\'eutilis\'ee par tous les g\'en\'erateurs de trajectoires qui partagent ce \cs{MarketEnvironment}.
Si la matrice n'est pas d\'efinie positive (notamment \`a cause de corr\'elations incoh\'erentes),
l'algorithme de Cholesky l\`eve une \cs{InvalidCorrelationMatrixException} imm\'ediatement
plut\^ot qu'au moment de la simulation.
\end{tcolorbox}

\subsubsection{\texttt{SimulatedPaths.cs} -- Objet de sortie des g\'en\'erateurs}

\cs{SimulatedPaths} est le type de retour de tous les \cs{IPathGenerator}. C'est le
seul objet qui circule entre le g\'en\'erateur et les payoffs.

\textbf{Structure :}
\begin{itemize}[leftmargin=1.5em]
  \item \cs{double[,,] Paths} -- tenseur de taille \cs{[nFactors, nPaths, nSteps+1]}.
        La dimension 0 correspond aux facteurs (actifs + taux FX), la dimension 1
        aux sc\'enarios, la dimension 2 au temps discr\'etis\'e ($S_0$ en index 0).
  \item \cs{AssetIndex[] AssetMapping} -- identifie chaque tranche de la dimension 0
        par son ticker, sa classe et sa devise.
  \item \cs{int NFactors, NPaths, NSteps} -- dimensions du tenseur.
\end{itemize}

\textbf{M\'ethodes :}
\begin{itemize}[leftmargin=1.5em]
  \item \cs{SimulatedPaths Extract(int[] factorIndices)} -- retourne un sous-tenseur
        avec uniquement les facteurs demand\'es. Utilis\'e par les payoffs mono-actif
        pour extraire leur tranche depuis une simulation conjointe de portefeuille.
  \item \cs{SimulatedPaths Extract(string[] tickers)} -- variante par nom de ticker.
  \item \cs{double[] GetPathForFactor(int factorIndex, int pathIndex)} -- acc\`es \`a une
        trajectoire individuelle.
  \item \cs{double[] GetTerminalValues(int factorIndex)} -- vecteur des spots terminaux
        $S_T^{(j)}$ pour tous les sc\'enarios. Utilis\'e par les payoffs europ\'eens purs.
\end{itemize}

\begin{tcolorbox}[note]
Le mapping \cs{AssetIndex[]} est fondamental. Il permet \`a un payoff basket de savoir
que la dimension 0 correspond au CAC40 et la dimension 1 au Nikkei, ind\'ependamment
de l'ordre dans lequel le simulateur a construit la matrice de corr\'elation.
Le portefeuille pricer (section 5.6) exploite ce m\'ecanisme pour distribuer les paths.
\end{tcolorbox}

\subsubsection{Instruments multi-actifs (inchang\'es de v1.2)}

\cs{BasketInstrument}, \cs{FxOptionInstrument}, \cs{QuantoInstrument},
\cs{FixingSchedule} : inchang\'es. Se reporter au walkthrough v1.2.

\subsection{Nouvelles interfaces}

\subsubsection{\texttt{IPathGenerator.cs}}

\begin{itemize}[leftmargin=1.5em]
  \item \cs{SimulatedPaths Generate(MarketEnvironment env, SimulationParameters params)} \\
        M\'ethode principale. Prend l'environnement de march\'e complet et les param\`etres
        de simulation, retourne le tenseur de trajectoires.
  \item \cs{int NFactors \{ get; \}} -- nombre de facteurs stochastiques g\'en\'er\'es.
  \item \cs{AssetIndex[] FactorMapping \{ get; \}} -- mapping des facteurs.
\end{itemize}

\subsubsection{\texttt{IPayoff.cs} -- Version 2.0}

L'interface \cs{IPayoff} est mise \`a jour pour prendre \cs{SimulatedPaths} en entr\'ee
(au lieu d'un \cs{double[]} mono-trajectoire).

\begin{itemize}[leftmargin=1.5em]
  \item \cs{double[] Compute(SimulatedPaths paths)} -- retourne un vecteur de payoffs
        non actualis\'es, un par sc\'enario. La vectorisation est native.
  \item \cs{string[] RequiredAssets \{ get; \}} -- tickers des actifs n\'ecessaires.
        Utilis\'e par le \cs{PortfolioPricer} pour construire la simulation conjointe.
  \item \cs{string Name \{ get; \}} -- lib\'ell\'e du produit.
\end{itemize}

\begin{tcolorbox}[note]
Le passage de \cs{double[]} \`a \cs{SimulatedPaths} est la rupture architecturale
principale entre v1.2 et v2.0. Les payoffs mono-actif s'adaptent trivalement :
ils appellent \cs{paths.Extract(RequiredAssets)} puis it\`erent sur les trajectoires.
Cette abstraction permet au \cs{PortfolioPricer} de passer une simulation conjointe
compl\`ete et \`a chaque payoff d'en extraire sa tranche.
\end{tcolorbox}

\subsubsection{\texttt{IVarianceReduction.cs}}

\begin{itemize}[leftmargin=1.5em]
  \item \cs{SimulatedPaths Apply(SimulatedPaths rawPaths, MarketEnvironment env)} \\
        Transforme ou enrichit les trajectoires brutes.
  \item \cs{string Name \{ get; \}}
\end{itemize}

\subsubsection{\texttt{IMonteCarloEngine.cs}}

\begin{itemize}[leftmargin=1.5em]
  \item \cs{Task<PricingResult> PriceAsync(IPathGenerator gen, IPayoff payoff,}\\
        \cs{MarketEnvironment env, SimulationParameters sim, CancellationToken ct)}
  \item \cs{Task<Greeks> ComputeGreeksAsync(IPathGenerator gen, IPayoff payoff,}\\
        \cs{MarketEnvironment env, SimulationParameters sim, CancellationToken ct)}
\end{itemize}

\subsubsection{\texttt{SimulationParameters.cs} -- Value object}

Record immuable portant les param\`etres de simulation :
\begin{itemize}[leftmargin=1.5em]
  \item \cs{int NumberOfPaths} (d\'efaut : 100\,000)
  \item \cs{int Steps} (d\'efaut : 252)
  \item \cs{double MaturityInYears}
  \item \cs{int? Seed} -- nullable pour reproductibilit\'e
\end{itemize}

\subsection{Interfaces inchang\'ees de v1.2}

\cs{IRepository<T>}, \cs{IPricer<TResult>}, \cs{IMarketDataFeed}, \cs{IRiskEngine},
\cs{IMultiAssetPricer<TResult>} : inchang\'ees.

\begin{tcolorbox}[deliverable]
Core v2.0 est termin\'e quand :
\begin{itemize}[leftmargin=1.5em]
  \item \cs{ProductType} compile avec tous les nouveaux membres et m\'ethodes d'extension coh\'erentes.
  \item \cs{GlobalCorrelationMatrix} l\`eve \cs{InvalidCorrelationMatrixException} sur une matrice
        non d\'efinie positive d\`es la construction.
  \item \cs{SimulatedPaths.Extract(string[])} retourne le bon sous-tenseur pour un panier de 3 actifs.
  \item \cs{MarketEnvironment.FromSingleAsset} permet aux tests v1.2 de compiler sans modification.
\end{itemize}
\end{tcolorbox}

\newpage
% ============================================================
\section{Cadre math\'ematique -- March\'e Multi-Devises sous Black-Scholes}
% ============================================================

\subsection{Description du march\'e}

On consid\`ere un march\'e international compos\'e de $1$ devise domestique et $\bar{n}$
devises \'etrang\`eres. On note $r^0$ le taux d'int\'er\^et domestique et $r^i$ le taux
d'int\'er\^et de la devise \'etrang\`ere $i \in \{1, \dots, \bar{n}\}$.

Le march\'e contient :
\begin{itemize}[leftmargin=1.5em]
  \item $n_0$ actifs risqu\'es domestiques $\bigl(S^{0,1}, \dots, S^{0,n_0}\bigr)$.
  \item Pour chaque devise $i$ : le taux de change $X^i$ (unit\'es de devise domestique
        pour une unit\'e de devise $i$) et $n_i$ actifs risqu\'es \'etrangers
        $\bigl(S^{i,1}, \dots, S^{i,n_i}\bigr)$.
\end{itemize}

Le nombre total de facteurs stochastiques est $N = \displaystyle\sum_{i=0}^{\bar{n}} n_i + \bar{n}$.

\subsection{Dynamiques sous la probabilit\'e historique $\mathbb{P}$}

\begin{align}
  \dd S^{0,\ell}_t &= S^{0,\ell}_t \bigl(\mu^{0,\ell}\,\dd t + \sigma^{0,\ell}\,\dd W^{0,\ell}_t\bigr),
  \quad \ell \in \{1,\dots,n_0\} \\
  \dd S^{i,\ell}_t &= S^{i,\ell}_t \bigl(\mu^{i,\ell}\,\dd t + \sigma^{i,\ell}\,\dd W^{i,\ell}_t\bigr),
  \quad i \in \{1,\dots,\bar{n}\},\; \ell \in \{1,\dots,n_i\} \\
  \dd X^i_t &= X^i_t \bigl(\mu^{i,X}\,\dd t + \sigma^{i,X}\,\dd W^{i,X}_t\bigr),
  \quad i \in \{1,\dots,\bar{n}\}
\end{align}

o\`u tous les mouvements browniens sont r\'eels et corr\'el\'es. La matrice de corr\'elation
$\rho \in \mathcal{M}_N(\mathbb{R})$ correspond au vecteur d'actifs $(\bar{S}, \bar{X})$
ordonn\'e comme suit :

\begin{enumerate}[leftmargin=1.5em]
  \item Les $n_0$ actifs risqu\'es domestiques.
  \item Pour $i \in \{1, \dots, \bar{n}\}$, les $n_i$ actifs risqu\'es dans la devise $i$.
  \item Pour $i \in \{1, \dots, \bar{n}\}$, le taux de change $X^i$.
\end{enumerate}

\textbf{Cet ordonnancement est la convention de \cs{GlobalCorrelationMatrix} et doit
\^etre respect\'e strictement par tous les g\'en\'erateurs de trajectoires.}

\subsection{Dynamiques sous la mesure risque-neutre domestique $\mathbb{Q}^0$}

Par le th\'eor\`eme de Girsanov et la condition d'absence d'arbitrage (parit\'e couverte
des taux d'int\'er\^et pour les FX), sous $\mathbb{Q}^0$ :

\begin{align}
  \dd S^{0,\ell}_t &= S^{0,\ell}_t \Bigl(r^0\,\dd t + \sigma^{0,\ell}\,\dd \widetilde{W}^{0,\ell}_t\Bigr)
  \label{eq:dom_asset} \\
  \dd X^i_t &= X^i_t \Bigl(\bigl(r^0 - r^i\bigr)\dd t + \sigma^{i,X}\,\dd \widetilde{W}^{i,X}_t\Bigr)
  \label{eq:fx_rate} \\
  \dd S^{i,\ell}_t &= S^{i,\ell}_t \Bigl(\bigl(r^i - \rho_{i\ell,\,iX}\,\sigma^{i,\ell}\,\sigma^{i,X}\bigr)
  \dd t + \sigma^{i,\ell}\,\dd \widetilde{W}^{i,\ell}_t\Bigr)
  \label{eq:foreign_asset}
\end{align}

o\`u $\rho_{i\ell,\,iX} = \text{Corr}(W^{i,\ell}, W^{i,X})$ est la corr\'elation entre
l'actif \'etranger $S^{i,\ell}$ et son taux de change $X^i$.

\begin{tcolorbox}[math]
\textbf{Ajustement quanto.} Le terme $-\rho_{i\ell,\,iX}\,\sigma^{i,\ell}\,\sigma^{i,X}$
dans la d\'erive de $S^{i,\ell}$ est l'ajustement quanto. Il compense le fait que, sous
$\mathbb{Q}^0$, un actif \'etranger pay\'e en devise domestique voit sa d\'erive modifi\'ee
par la corr\'elation entre le rendement de l'actif et le rendement du change. Si
$\rho_{i\ell,\,iX} > 0$ (actif et change corr\'el\'es positivement), la d\'erive est r\'eduite.
Pour une option quanto pure ($X^i$ fix\'e contractuellement), seule l'\'equation
\eqref{eq:foreign_asset} intervient dans la simulation.
\end{tcolorbox}

\subsection{Discr\'etisation (sch\'ema log-Euler exact sous GBM)}

Le sch\'ema log-Euler est exact pour le GBM (pas de biais de discr\'etisation) :

\begin{align}
  S^{0,\ell}_{t+\Delta t} &= S^{0,\ell}_t \exp\!\Bigl[\Bigl(r^0 - \tfrac{(\sigma^{0,\ell})^2}{2}\Bigr)\Delta t
    + \sigma^{0,\ell}\sqrt{\Delta t}\,\widetilde{Z}^{0,\ell}\Bigr] \\
  X^i_{t+\Delta t} &= X^i_t \exp\!\Bigl[\Bigl(r^0 - r^i - \tfrac{(\sigma^{i,X})^2}{2}\Bigr)\Delta t
    + \sigma^{i,X}\sqrt{\Delta t}\,\widetilde{Z}^{i,X}\Bigr] \\
  S^{i,\ell}_{t+\Delta t} &= S^{i,\ell}_t \exp\!\Bigl[\Bigl(r^i - \rho_{i\ell,iX}\sigma^{i,\ell}\sigma^{i,X}
    - \tfrac{(\sigma^{i,\ell})^2}{2}\Bigr)\Delta t + \sigma^{i,\ell}\sqrt{\Delta t}\,\widetilde{Z}^{i,\ell}\Bigr]
\end{align}

o\`u $\widetilde{\mathbf{Z}} = L \boldsymbol{\xi}$ avec $\boldsymbol{\xi} \sim \mathcal{N}(0, I_N)$
et $L$ le facteur de Cholesky de $\rho$ (tel que $LL^\top = \rho$).

\subsection{Structure de payoff g\'en\'eral}

Pour des dates de constatation $0 \le t_1 < \cdots < t_M = T$, l'outil peut pricer
tout flux pay\'e en $T$ de la forme :

\[
  \psi\!\bigl(\bar{S}_{t_1}, \dots, \bar{S}_{t_M},\, \bar{X}_{t_1}, \dots, \bar{X}_{t_M}\bigr)
\]

Le prix \`a $t=0$ sous $\mathbb{Q}^0$ est :

\[
  P_0 = e^{-r^0 T}\,\mathbb{E}^{\mathbb{Q}^0}\!\bigl[\psi\bigl(\bar{S}_{t_1}, \dots, \bar{S}_{t_M},
  \bar{X}_{t_1}, \dots, \bar{X}_{t_M}\bigr)\bigr]
\]

Estim\'e par Monte Carlo : $\hat{P}_0 = e^{-r^0 T} \cdot \frac{1}{n}\sum_{j=1}^n
\psi^{(j)}$ avec erreur standard $\widehat{\text{SE}} = e^{-r^0 T}\cdot\hat{\sigma}_\psi/\sqrt{n}$
et intervalle de confiance \`a 95\% : $\hat{P}_0 \pm 1.96\,\widehat{\text{SE}}$.

\newpage
% ============================================================
\section{Module Pricing v2.0 -- \texttt{TradingDashboard.Pricing}}
% ============================================================

\subsection{Structure des dossiers}

\begin{itemize}[leftmargin=1.5em]
  \item \textbf{BlackScholes/} -- pricer analytique vanilles (inchang\'e de v1.2).
  \item \textbf{PathGenerators/} -- g\'en\'erateurs de trajectoires (nouveau en v2.0).
  \item \textbf{Payoffs/} -- tous les payoffs (\'etendu en v2.0).
  \item \textbf{VarianceReduction/} -- r\'eductions de variance (nouveau en v2.0).
  \item \textbf{Engine/} -- moteur MC et pricer de portefeuille (nouveau en v2.0).
  \item \textbf{Greeks/} -- calculateur de sensibilit\'es (inchang\'e).
  \item \textbf{Strategy/} -- factories et strat\'egies (\'etendu en v2.0).
\end{itemize}

% ============================================================
\subsection{BlackScholes/ -- Inchang\'e de v1.2}

\cs{BlackScholesFormulas.cs}, \cs{VanillaOptionPricer.cs}, \cs{GreeksCalculator.cs},
\cs{PricingResult.cs} : identiques \`a la v1.2. Se reporter \`a la section correspondante.
La m\'ethode \cs{VanillaOptionPricer} accepte d\'esormais un \cs{MarketEnvironment}
en plus du \cs{MarketSnapshot} (via \cs{FromSingleAsset}).

% ============================================================
\subsection{PathGenerators/}

\subsubsection{\texttt{GbmPathGenerator.cs} -- Mono-actif (refactor de GBMPathSimulator)}

Impl\'emente \cs{IPathGenerator}. Encapsule l'ancien \cs{GBMPathSimulator} de v1.2
et retourne un \cs{SimulatedPaths} de dimensions $[1, N_\text{paths}, N_\text{steps}+1]$.

Aucune logique math\'ematique chang\'ee. Seule l'interface de retour est mise \`a jour.
L'\cs{AssetMapping} contient un unique \'el\'ement avec le ticker du sous-jacent.

\textbf{R\'etrocompatibilit\'e :} une m\'ethode \cs{SimulatePath(Random rng) : double[]} est
conserv\'ee pour les tests unitaires existants.

\subsubsection{\texttt{MultiAssetGbmPathGenerator.cs} -- Multi-actifs domestiques}

\textbf{Cas couvert :} $\bar{n} = 0$ (pas de devise \'etrang\`ere). Tous les actifs sont
dans la devise domestique. Taille du vecteur facteur : $N = n_0$.

\textbf{Algorithme :}
\begin{enumerate}[leftmargin=1.5em]
  \item Lire la sous-matrice de corr\'elation $\rho_0 \in \mathcal{M}_{n_0}(\mathbb{R})$
        depuis \cs{env.Correlation.Sub(domesticIndices)}.
  \item R\'ecup\'erer le facteur de Cholesky $L_0$ (pr\'e-calcul\'e dans \cs{GlobalCorrelationMatrix}).
  \item Pour chaque pas de temps $\Delta t$ et chaque sc\'enario $j$ :
        \begin{enumerate}[leftmargin=1.5em]
          \item G\'en\'erer $\boldsymbol{\xi} \sim \mathcal{N}(0, I_{n_0})$.
          \item Calculer $\widetilde{\mathbf{Z}} = L_0\,\boldsymbol{\xi}$.
          \item Appliquer l'\'equation \eqref{eq:dom_asset} pour chaque actif $\ell$.
        \end{enumerate}
  \item Retourner \cs{SimulatedPaths} avec tenseur $[n_0, N_\text{paths}, N_\text{steps}+1]$.
\end{enumerate}

\begin{tcolorbox}[constraint]
\textbf{Pas de r\'eallocation par sc\'enario.} Allouer le tenseur de sortie une seule
fois avant la boucle principale. \cs{Parallel.For} avec \cs{ThreadLocal<Random>}
sur la dimension sc\'enario. Chaque thread g\`ere un sous-ensemble de sc\'enarios.
\end{tcolorbox}

\subsubsection{\texttt{MultiCurrencyGbmPathGenerator.cs} -- Multi-devises (cas g\'en\'eral)}

Impl\'emente le mod\`ele complet de la section~3. Prend en charge $\bar{n} \ge 1$ devises
\'etrang\`eres avec leurs actifs et taux FX.

\textbf{Taille du vecteur facteur :} $N = \sum_{i=0}^{\bar{n}} n_i + \bar{n}$.

\textbf{Algorithme :}
\begin{enumerate}[leftmargin=1.5em]
  \item V\'erifier que l'\cs{AssetMapping} de \cs{env.Correlation} respecte l'ordonnancement
        de la section~3.2 (domestiques, puis \'etrangers par devise, puis FX).
  \item R\'ecup\'erer le facteur de Cholesky $L \in \mathcal{M}_N(\mathbb{R})$ depuis
        \cs{env.Correlation.CholeskyFactor}.
  \item Pour chaque sc\'enario $j$ et chaque pas de temps $\Delta t$ :
        \begin{enumerate}[leftmargin=1.5em]
          \item G\'en\'erer $\boldsymbol{\xi} \sim \mathcal{N}(0, I_N)$.
          \item Calculer $\widetilde{\mathbf{Z}} = L\,\boldsymbol{\xi}$.
          \item Appliquer \eqref{eq:dom_asset} aux actifs domestiques.
          \item Appliquer \eqref{eq:fx_rate} aux taux de change $X^i$.
          \item Appliquer \eqref{eq:foreign_asset} aux actifs \'etrangers, en lisant
                $\rho_{i\ell,iX}$ dans la matrice de corr\'elation globale.
        \end{enumerate}
  \item Retourner \cs{SimulatedPaths} avec tenseur $[N, N_\text{paths}, N_\text{steps}+1]$.
\end{enumerate}

\textbf{Cas d\'eg\'en\'er\'es (test d'architecture) :}
\begin{itemize}[leftmargin=1.5em]
  \item $\bar{n}=0$, $n_0=1$ : doit produire des trajectoires identiques (en loi) au \cs{GbmPathGenerator}.
  \item $\bar{n}=1$, $n_1=0$ : simulation d'un seul taux FX, v\'erifier que la d\'erive est $r^0 - r^1$.
\end{itemize}

\begin{tcolorbox}[todo]
\textbf{HestonPathGenerator, LocalVolPathGenerator, LSVPathGenerator.} Ces g\'en\'erateurs
impl\'ementeront la m\^eme interface \cs{IPathGenerator}. Ils seront d\'evelopp\'es dans
des it\'erations ult\'erieures apr\`es calibration de la surface de volatilit\'e. L'architecture
de la v2.0 garantit qu'aucun payoff ni aucun moteur MC ne devra \^etre modifi\'e lors
de leur introduction.
\end{tcolorbox}

% ============================================================
\subsection{Payoffs/}

Tous les payoffs impl\'ementent \cs{IPayoff}. La m\'ethode \cs{Compute(SimulatedPaths paths)}
commence syst\'ematiquement par \cs{paths.Extract(RequiredAssets)} pour isoler la
tranche pertinente.

\subsubsection{Payoffs mono-actif existants (migration v1.2 $\to$ v2.0)}

\cs{AsianCallPayoff}, \cs{AsianPutPayoff}, \cs{LookbackFixedStrikeCallPayoff},
\cs{LookbackFixedStrikePutPayoff} : logique de calcul inchang\'ee, seule la signature
de \cs{Compute} est mise \`a jour pour prendre \cs{SimulatedPaths}.

\subsubsection{\texttt{Payoffs/FloatingLookbackCallPayoff.cs} et \texttt{FloatingLookbackPutPayoff.cs}}

\begin{tcolorbox}[math]
\textbf{Lookback \`a strike flottant.}
\[
  \text{FloatingLookbackCall} = \max\!\Bigl(S_T - \min_{t \in [0,T]} S_t,\; 0\Bigr)
  = S_T - \min_{t} S_t
\]
\[
  \text{FloatingLookbackPut} = \max\!\Bigl(\max_{t \in [0,T]} S_t - S_T,\; 0\Bigr)
  = \max_{t} S_t - S_T
\]
Les deux expressions de droite sont exactes car le floating lookback call (resp. put)
est toujours $\ge 0$ par construction. Des formules ferm\'ees existent sous GBM
(Goldman-Sosin-Gatto 1979) et servent de validation des r\'esultats Monte Carlo.
\end{tcolorbox}

Impl\'ementation : \cs{path.Min()} et \cs{path.Max()} sur la trajectoire compl\`ete.

\subsubsection{\texttt{Payoffs/BarrierPayoffs.cs} -- Famille barri\`ere}

\begin{tcolorbox}[math]
\textbf{Payoffs barri\`eres (type call, g\'en\'eralisable au put).}

Soit $B$ le niveau de barri\`ere, $\mathbf{1}_A$ l'indicatrice de l'\'ev\'enement $A$.

\vspace{0.3em}
\textit{Down-and-Out Call :}
\[
  \text{DaO Call} = \max(S_T - K, 0)\cdot\mathbf{1}\!\left\{\min_{t \in [0,T]} S_t > B\right\}
\]
\textit{Down-and-In Call :}
\[
  \text{DaI Call} = \max(S_T - K, 0)\cdot\mathbf{1}\!\left\{\min_{t \in [0,T]} S_t \le B\right\}
\]
\textit{Up-and-Out Call :}
\[
  \text{UaO Call} = \max(S_T - K, 0)\cdot\mathbf{1}\!\left\{\max_{t \in [0,T]} S_t < B\right\}
\]
\textit{Double Barrier Call} (deux barri\`eres $B_{\text{low}} < B_{\text{high}}$) :
\[
  \text{DB Call} = \max(S_T - K, 0)\cdot\mathbf{1}\!\left\{B_{\text{low}} < \min_t S_t \text{ et } \max_t S_t < B_{\text{high}}\right\}
\]
\end{tcolorbox}

\textbf{Correction de Broadie-Glasserman-Kou (BGK) pour la discr\'etisation.}
En monitoring discret avec $M$ pas de temps, les barri\`eres sont syst\'ematiquement
biais\'ees car on rate les extrema intra-p\'eriode. La correction BGK d\'ecale la barri\`ere
de $\pm\beta\sigma\sqrt{\Delta t}$ avec $\beta \approx 0.5826$ :

\[
  B_{\text{eff}}^{\text{down}} = B \cdot e^{-\beta\sigma\sqrt{\Delta t}},\quad
  B_{\text{eff}}^{\text{up}} = B \cdot e^{+\beta\sigma\sqrt{\Delta t}}
\]

\textbf{Impl\'ementation :} une classe abstraite \cs{BarrierPayoffBase} factorise le
calcul de franchissement et l'application de la correction BGK. Les 10 types concrets
h\'eritent de cette base et red\'efinissent uniquement la condition de franchissement.

\textbf{Param\`etres constructeur (exemple \cs{DownAndOutCallPayoff}) :}
\cs{(double strike, double barrier, double vol, double maturity, int steps, bool applyBgk = true)}

\begin{tcolorbox}[note]
La correction BGK est activ\'ee par d\'efaut. En tests de validation contre les formules
ferm\'ees BS (Rubinstein-Reiner), la d\'esactiver (\cs{applyBgk = false}) et augmenter
\cs{steps} jusqu'\`a convergence. L'\'ecart entre le r\'esultat avec et sans BGK mesure
le biais de discr\'etisation r\'esiduel.
\end{tcolorbox}

\subsubsection{\texttt{Payoffs/DigitalPayoffs.cs}}

\begin{tcolorbox}[math]
\textbf{Digitales.}

\textit{Cash-or-Nothing Call} (paiement $N$ si $S_T > K$) :
\[
  \text{CaN Call} = N \cdot \mathbf{1}\{S_T > K\}
\]
\textit{Asset-or-Nothing Call} (paiement $S_T$ si $S_T > K$) :
\[
  \text{AaN Call} = S_T \cdot \mathbf{1}\{S_T > K\}
\]
Formules ferm\'ees sous BS : $\text{CaN Call} = N \cdot e^{-rT} N(d_2)$,
$\text{AaN Call} = S_0\, N(d_1)$. Servent de validation Monte Carlo.
\end{tcolorbox}

\textbf{Note de stabilit\'e num\'erique.} Les digitales ont un vega tr\`es peak\'e pr\`es
du strike \`a maturit\'e. La m\'ethode brute Monte Carlo convergera lentement. Une
am\'elioration possible (non exig\'ee en v2.0) est le lissage du payoff indicateur
(\textit{fuzzy boundary}).

\subsubsection{\texttt{Payoffs/VarianceSwapPayoff.cs}}

\begin{tcolorbox}[math]
\textbf{Variance Swap.} Notionnel $N_\text{var}$, strike de variance $K_\text{var}$.

La variance r\'ealis\'ee discr\`ete sur $M$ retours log est :
\[
  \sigma^2_\text{r\'ealis\'ee} = \frac{252}{M} \sum_{j=1}^{M} \left[\ln\frac{S_{t_j}}{S_{t_{j-1}}}\right]^2
\]

Payoff du vendeur de variance swap (paiement en $T$) :
\[
  \text{VarianceSwap} = N_\text{var} \cdot \bigl(\sigma^2_\text{r\'ealis\'ee} - K_\text{var}\bigr)
\]

Sous BS, $\mathbb{E}^{\mathbb{Q}}[\sigma^2_\text{r\'ealis\'ee}] = \sigma^2_\text{implicite}$
exactement en monitoring continu. Le pricer Monte Carlo permet de quantifier le biais
de monitoring discret.
\end{tcolorbox}

\subsubsection{\texttt{Payoffs/VolatilitySwapPayoff.cs}}

\begin{tcolorbox}[math]
\textbf{Volatility Swap.} Payoff analogique mais sur la volatilit\'e r\'ealis\'ee :
\[
  \text{VolatilitySwap} = N_\text{vol} \cdot \bigl(\sigma_\text{r\'ealis\'ee} - K_\text{vol}\bigr),
  \quad \sigma_\text{r\'ealis\'ee} = \sqrt{\sigma^2_\text{r\'ealis\'ee}}
\]
Pas de r\'eplication statique exacte m\^eme sous BS. Monte Carlo est la m\'ethode de
r\'ef\'erence. $\mathbb{E}[\sigma_\text{r\'ealis\'ee}] < \sqrt{\mathbb{E}[\sigma^2_\text{r\'ealis\'ee}]}$
par l'in\'egalit\'e de Jensen, d'o\`u la \textit{convexity adjustment} entre variance swap
et volatility swap.
\end{tcolorbox}

\subsubsection{\texttt{Payoffs/CliquetPayoff.cs}}

\begin{tcolorbox}[math]
\textbf{Cliquet (\textit{resetted-strike}).} Dates de reset $0 = T_0 < T_1 < \cdots < T_M = T$,
cap $c \ge 0$, floor $f \le 0$.

Le retour local plafonn\'e-planch\'e sur la p\'eriode $[T_{k-1}, T_k]$ est :
\[
  r_k = \max\!\Bigl(f,\;\min\!\Bigl(c,\; \frac{S_{T_k} - S_{T_{k-1}}}{S_{T_{k-1}}}\Bigr)\Bigr)
\]
Payoff global :
\[
  \text{Cliquet} = \max\!\Bigl(0,\; \sum_{k=1}^M r_k\Bigr)
\]
\textbf{Sensibilit\'e au mod\`ele.} Le cliquet d\'epend de la \textit{forward volatility}
(vol implicite pour des options forward-start), non directement observable. Deux mod\`eles
calibr\'es identiquement \`a la m\^eme surface de vol statique peuvent produire des prix
de cliquet tr\`es diff\'erents selon leur dynamique de vol future.
\end{tcolorbox}

\subsubsection{\texttt{Payoffs/MultiAsset/BasketPayoff.cs}}

\begin{tcolorbox}[math]
\textbf{Basket Call.} Poids $w_i \ge 0$, $\sum_i w_i = 1$, actifs $S^1, \dots, S^n$.
\[
  \text{BasketCall} = \max\!\Bigl(\sum_{i=1}^n w_i S^i_T - K,\; 0\Bigr)
\]
Pas de formule ferm\'ee pour la somme pond\'er\'ee log-normale. Monte Carlo est exact.
\end{tcolorbox}

\cs{RequiredAssets} retourne tous les tickers du panier. \cs{Compute} extrait chaque
tranche du tenseur, construit la moyenne pond\'er\'ee et applique le max.

\subsubsection{\texttt{Payoffs/MultiAsset/WorstOfPayoff.cs} et \texttt{BestOfPayoff.cs}}

\begin{tcolorbox}[math]
\textbf{Worst-of / Best-of Call.} Actifs $S^1, \dots, S^n$ exprim\'es en rendement.
\[
  \text{WorstOfCall} = \max\!\Bigl(\min_{i} \frac{S^i_T}{S^i_0} \cdot S^{\text{ref}}_0 - K,\; 0\Bigr)
\]
\[
  \text{BestOfCall} = \max\!\Bigl(\max_{i} S^i_T - K,\; 0\Bigr)
\]
Le worst-of est structurellement plus cher en termes de prime \`a vendre mais moins
cher \`a acheter qu'un basket, car le payoff est minor\'e par le pire performer. Il est
le constituant de base de tous les autocalls.
\end{tcolorbox}

\subsubsection{\texttt{Payoffs/MultiAsset/SpreadPayoff.cs}}

\begin{tcolorbox}[math]
\textbf{Spread Option.} Deux actifs $S^1$ et $S^2$.
\[
  \text{SpreadOption} = \max(S^1_T - S^2_T - K,\; 0)
\]
Cas particulier $K = 0$ : option d'\'echange de Margrabe, avec formule ferm\'ee
$\text{Margrabe}(S^1, S^2, \sigma^1, \sigma^2, \rho_{12}, T)$ servant de validation.
\end{tcolorbox}

\subsubsection{\texttt{Payoffs/MultiAsset/QuantoPayoff.cs}}

\begin{tcolorbox}[math]
\textbf{Quanto Call.} Actif \'etranger $S^{i,\ell}$, taux FX $X^i$ fix\'e contractuellement
au niveau $\bar{X}$ (d\'eterministe). Paiement en devise domestique.
\[
  \text{QuantoCall} = \bar{X} \cdot \max(S^{i,\ell}_T - K,\; 0)
\]
La d\'erive de $S^{i,\ell}$ sous $\mathbb{Q}^0$ est $r^i - \rho_{i\ell,iX}\sigma^{i,\ell}\sigma^{i,X}$
(\'equation~\eqref{eq:foreign_asset}) : l'ajustement quanto est intrins\`eque \`a la
simulation. La simulation du taux FX $X^i$ n'est pas n\'ecessaire pour ce payoff
(il est fix\'e contractuellement), mais $\rho_{i\ell,iX}$ intervient dans la d\'erive.
\end{tcolorbox}

\subsubsection{Payoffs structured products (placeholders)}

\begin{tcolorbox}[todo]
\textbf{AutocallAthenaPayoff, AutocallPhoenixPayoff, ReverseConvertiblePayoff.}

Ces payoffs requi\`erent un moteur MC sous LSV (Local Stochastic Volatility) pour
\^etre \'economiquement pertinents. Sous BS, le pricer produit une r\'ef\'erence valide
en structure (test de signe, mono-convexity) mais pas un prix de march\'e.

L'interface \cs{IPayoff} est d\'ej\`a adapt\'ee. Ces classes seront impl\'ement\'ees lors
de l'introduction du \cs{LSVPathGenerator}. Activer des \cs{NotImplementedException}
pour l'instant.
\end{tcolorbox}

% ============================================================
\subsection{VarianceReduction/}

\subsubsection{\texttt{BrownianBridgeVarianceReduction.cs}}

Impl\'emente \cs{IVarianceReduction}. Applicable aux payoffs \`a barri\`ere pour
corriger le biais de discr\'etisation via la probabilit\'e de franchissement intra-p\'eriode.

\begin{tcolorbox}[math]
\textbf{Pont brownien pour les barrières.} Soit $S_k = S_{t_k}$ et $S_{k+1} = S_{t_{k+1}}$
les spots simulés aux deux extrémités d'un intervalle $[t_k, t_{k+1}]$. La probabilit\'e
que le processus GBM franchisse une barri\`ere inf\'erieure $B$ entre $t_k$ et $t_{k+1}$
conditionnellement aux valeurs aux bornes est :

\[
  p_k = \exp\!\left(-\frac{2\,\ln(S_k/B)\,\ln(S_{k+1}/B)}{\sigma^2\Delta t}\right)
\]

Le payoff knock-out est multipli\'e par $\prod_k (1 - p_k)$ au lieu d'utiliser
l'indicatrice exacte (qui n\'ecessiterait des pas tr\`es fins). Cette correction
\'elimine le biais de discr\'etisation \`a l'ordre dominant.
\end{tcolorbox}

\subsubsection{\texttt{GeometricAsianControlVariate.cs}}

Impl\'emente \cs{IVarianceReduction}. G\'en\`ere simultan\'ement la moyenne g\'eom\'etrique
et arithm\'etique sur le m\^eme path.

\begin{tcolorbox}[math]
\textbf{Variable de contr\^ole g\'eom\'etrique.} Soit $\bar{S}_\text{arith}$ la moyenne
arithm\'etique et $\bar{S}_\text{g\'eom}$ la moyenne g\'eom\'etrique sur les m\^emes
observations. La moyenne g\'eom\'etrique a une formule ferm\'ee sous BS (Kemna-Vorst).
On exploite la forte corr\'elation entre $\psi_\text{arith}$ et $\psi_\text{g\'eom}$.

Estimateur corrig\'e :
\[
  \hat{P}_\text{arith} = \bar{\psi}_\text{arith}
    - \hat{\beta}\,(\bar{\psi}_\text{g\'eom} - P_\text{g\'eom}^{\text{BS}})
\]
o\`u $\hat{\beta}$ est estim\'e par r\'egression OLS sur les sc\'enarios et
$P_\text{g\'eom}^{\text{BS}}$ est le prix analytique exact de l'option g\'eom\'etrique.
\end{tcolorbox}

\subsubsection{\texttt{AntitheticVarianceReduction.cs}}

Pour chaque sc\'enario $\boldsymbol{\xi}^{(j)}$, g\'en\`ere le sc\'enario antith\'etique
$-\boldsymbol{\xi}^{(j)}$ et retourne la moyenne des deux payoffs. R\'eduit la variance
sans biais. Applicable \`a tous les payoffs.

% ============================================================
\subsection{Engine/}

\subsubsection{\texttt{MonteCarloEngine.cs} -- Moteur g\'en\'eral}

Impl\'emente \cs{IMonteCarloEngine}. Orchestre le g\'en\'erateur, la r\'eduction de
variance et le payoff.

\textbf{Algorithme \cs{PriceAsync} :}
\begin{enumerate}[leftmargin=1.5em]
  \item \cs{gen.Generate(env, sim)} -- tenseur de trajectoires.
  \item Appliquer \cs{IVarianceReduction.Apply(paths, env)} si pr\'esente.
  \item \cs{payoff.Compute(paths)} -- vecteur de payoffs bruts $\psi^{(j)}$.
  \item Facteur d'actualisation : $D = e^{-r^0 T}$.
  \item $\hat{P} = D \cdot \frac{1}{n}\sum_j \psi^{(j)}$.
  \item $\hat{\sigma}_\psi^2 = \frac{1}{n-1}\sum_j (\psi^{(j)} - \bar{\psi})^2$.
  \item $\text{SE} = D \cdot \hat{\sigma}_\psi / \sqrt{n}$, $\text{CI}_{95} = 1.96 \cdot \text{SE}$.
  \item Retourner \cs{PricingResult} avec prix, SE, CI, m\'ethode, horodatage.
\end{enumerate}

\textbf{M\'ethode \cs{ComputeGreeksAsync} -- diff\'erences finies centrales :}

Tous les bumps sont lanc\'es en parall\`ele via \cs{Task.WhenAll}. Le g\'en\'erateur
utilise la m\^eme seed pour garantir la comparabilit\'e des simulations bump\'ees
et non bump\'ees (common random numbers).

\begin{center}
\begin{tabularx}{\textwidth}{l l X}
\toprule
\textbf{Greek} & \textbf{Formule} & \textbf{Bump} \\
\midrule
Delta & $\frac{P(S_0+h) - P(S_0-h)}{2h}$ & $h = 0.01 \times S_0$ \\
Gamma & $\frac{P(S_0+h) - 2P(S_0) + P(S_0-h)}{h^2}$ & m\^eme $h$ \\
Vega  & $\frac{P(\sigma+\delta\sigma) - P(\sigma-\delta\sigma)}{2\,\delta\sigma}$ & $\delta\sigma = 0.01$ \\
Theta & $\frac{P(T-\delta T) - P(T)}{365\,\delta T}$ & $\delta T = 1/365$ \\
Rho   & $\frac{P(r+\delta r) - P(r-\delta r)}{2\,\delta r}$ & $\delta r = 0.0001$ \\
\bottomrule
\end{tabularx}
\end{center}

\begin{tcolorbox}[constraint]
\textbf{Common random numbers pour les Greeks.} Si \cs{Seed} est renseign\'e dans
\cs{SimulationParameters}, les simulations bump\'ees utilisent la m\^eme seed que la
simulation de base. Cela r\'eduit drastiquement la variance des estimateurs de Greeks
par diff\'erences finies (la variabilit\'e MC se compense entre les deux appels bump\'es).
Sans cette technique, la variance du Delta MC peut \^etre plus grande que le signal.
\end{tcolorbox}

\subsubsection{\texttt{PortfolioPricer.cs} -- Simulation conjointe}

C'est la pi\`ece ma\^itresse de l'architecture v2.0 pour le pricing de portefeuille.

\textbf{Probl\`eme.} Si l'on price chaque position ind\'ependamment, les corr\'elations
entre sous-jacents sont ignor\'ees. Un portefeuille long SPX / short Eurostoxx n'est
pas la somme de deux positions isol\'ees~: leur corr\'elation d\'etermine le risque r\'eel.

\textbf{Algorithme :}
\begin{enumerate}[leftmargin=1.5em]
  \item Collecter \cs{RequiredAssets} de chaque instrument du portefeuille.
  \item D\'eduire l'univers global $\mathcal{U}$ (union sans doublon).
  \item Extraire la sous-matrice de corr\'elation $\rho_\mathcal{U}$ depuis
        \cs{env.Correlation.Sub(indices($\mathcal{U}$))}.
  \item Lancer \textbf{une unique simulation conjointe} avec
        \cs{MultiCurrencyGbmPathGenerator} sur $\mathcal{U}$.
  \item Pour chaque instrument $k$ : extraire sa tranche via
        \cs{paths.Extract(instrument.RequiredAssets)}.
  \item Calculer \cs{payoff.Compute(subPaths)} pour chaque position.
  \item PnL du portefeuille : $\text{PnL}^{(j)} = \sum_k q_k \cdot (D \cdot \psi_k^{(j)} - P_k^{(0)})$.
  \item Retourner le vecteur de PnL et le prix du portefeuille.
\end{enumerate}

\textbf{M\'ethodes :}
\begin{itemize}[leftmargin=1.5em]
  \item \cs{Task<PortfolioPricingResult> PriceAsync(IReadOnlyList<Position> positions,}\\
        \cs{MarketEnvironment env, SimulationParameters sim, CancellationToken ct)}
  \item \cs{Task<double[]> SimulatePortfolioPnLAsync(...)} -- retourne le vecteur
        brut des PnL sc\'enarios pour VaR/CVaR.
\end{itemize}

% ============================================================
\subsection{Strategy/ -- Factories et dispatch}

\subsubsection{\texttt{PricerFactory.cs} -- Version 2.0}

\begin{center}
\begin{tabularx}{\textwidth}{l l}
\toprule
\textbf{Condition} & \textbf{Strat\'egie retourn\'ee} \\
\midrule
\cs{IsVanilla()} & \cs{BlackScholesPricingStrategy} \\
\cs{IsExotic()} (mono-actif) & \cs{MonteCarloStrategy} avec \cs{GbmPathGenerator} \\
\cs{IsFx()} & \cs{MonteCarloStrategy} avec \cs{MultiCurrencyGbmPathGenerator} \\
\cs{IsMultiAsset()} & \cs{MonteCarloStrategy} avec \cs{MultiAssetGbmPathGenerator} \\
\cs{IsStructured()} & \cs{NotImplementedException} (placeholder) \\
\bottomrule
\end{tabularx}
\end{center}

\textbf{Invariant :} aucun \cs{if}/\cs{switch} sur \cs{ProductType} en dehors de
\cs{PricerFactory} et \cs{PayoffFactory}. Le dispatch est centralis\'e.

\begin{tcolorbox}[deliverable]
Pricing v2.0 est termin\'e quand :
\begin{itemize}[leftmargin=1.5em]
  \item \cs{MultiAssetGbmPathGenerator} : les paths de deux actifs corr\'el\'es \`a $\rho=0.8$
        ont une corr\'elation empirique dans $[0.79, 0.81]$ sur $10^5$ sc\'enarios.
  \item \cs{DownAndOutCallPayoff} avec \cs{applyBgk = false} converge vers la formule
        Rubinstein-Reiner \`a $\pm 0.5\%$ pour \cs{steps = 1000}.
  \item \cs{FloatingLookbackCallPayoff} : prix MC $\in \pm 1\%$ de Goldman-Sosin-Gatto
        pour $10^5$ sc\'enarios, \cs{steps = 252}.
  \item \cs{VarianceSwapPayoff} : $\mathbb{E}[\sigma^2_\text{r\'ealis\'ee}] \approx \sigma^2_\text{impl.}$
        \`a $\pm 0.1\%$ pour monitoring continu (grand $M$).
  \item \cs{PortfolioPricer} sur un long call SPX + short call Eurostoxx corr\'el\'es :
        le prix du portefeuille diff\`ere du somme des prix individuels (preuve que la
        simulation est bien conjointe).
\end{itemize}
\end{tcolorbox}

\newpage
% ============================================================
\section{Module Risk v2.0 -- \texttt{TradingDashboard.Risk}}
% ============================================================

\begin{tcolorbox}[note]
\textbf{Distinction fondamentale Pricing MC vs Risk MC.}
Le Monte Carlo de \cs{.Pricing} price un \textit{instrument individuel} : il calcule
$E[e^{-rT}\,\text{Payoff}(S_T)]$. Le Monte Carlo de \cs{.Risk} simule un
\textit{portefeuille} : il distribue les pertes sur $H$ jours et en d\'eduit la VaR et la CVaR.
Ce sont deux usages distincts du m\^eme \cs{GBMPathSimulator} -- d'o\`u son emplacement
dans \cs{.Pricing} accessible par les deux modules.
\end{tcolorbox}

\subsection{MonteCarlo/ -- Simulation conjointe de portefeuille}

\subsubsection{\texttt{PortfolioSimulator.cs}}

\textbf{R\^ole :} simuler la distribution de PnL d'un portefeuille sur un horizon
$H$ (typiquement 1 jour ou 10 jours) en g\'en\'erant des trajectoires GBM pour chaque
sous-jacent du portefeuille. Retourne la distribution empirique des pertes, la VaR
et la CVaR.

\textbf{Mise \`a jour critique par rapport \`a v1.2.} La v1.2 simulait chaque
sous-jacent ind\'ependamment, d\'etruisant les corr\'elations. La v2.0 utilise
\cs{PortfolioPricer.SimulatePortfolioPnLAsync} pour une simulation conjointe.

\textbf{Constructeur :}
\begin{itemize}[leftmargin=1.5em]
  \item \cs{PortfolioSimulator(GBMPathSimulator simulator, IPricer<PricingResult> pricer,}\\
        \cs{int numberOfScenarios = 10\_000)}
  \item \cs{GBMPathSimulator} inject\'e depuis \cs{.Pricing} -- pas de \cs{new}
        interne. C'est la r\'eutilisation inter-modules \'evoqu\'ee dans l'architecture.
\end{itemize}

\textbf{Algorithme VaR/CVaR :}
\begin{enumerate}[leftmargin=1.5em]
  \item Appeler \cs{portfolioPricer.SimulatePortfolioPnLAsync(positions, env,}\\
        \cs{sim with \{ MaturityInYears = horizon/252.0 \}, ct)}.
  \item Trier le vecteur de PnL $(\text{PnL}^{(1)}, \dots, \text{PnL}^{(n)})$ par ordre croissant.
  \item $\text{VaR}_{1-\alpha} = -\text{PnL}_{[\alpha n]}$ (l'\`eme percentile inf\'erieur).
  \item $\text{CVaR}_{1-\alpha} = -\dfrac{1}{\lfloor\alpha n\rfloor}\sum_{j \le \alpha n}\text{PnL}^{(j)}$.
  \item V\'erifier l'invariant : $\text{CVaR} \ge \text{VaR}$ (propri\'et\'e de coh\'erence).
\end{enumerate}

\textbf{R\'etrocompatibilit\'e :} pour un portefeuille mono-actif, \cs{PortfolioSimulator}
se comporte comme en v1.2 gr\^ace \`a \cs{MarketEnvironment.FromSingleAsset}.

\textbf{M\'ethodes principales :}
\begin{itemize}[leftmargin=1.5em]
  \item \cs{Task<VaRResult> ComputeVaRAsync(IReadOnlyList<Position> positions,}\\
        \cs{MarketSnapshot snapshot, int horizon = 1,}\\
        \cs{double confidenceLevel = 0.99, CancellationToken ct = default)}
  \item \cs{Task<CVaRResult> ComputeCVaRAsync(...)} -- m\^eme signature, calcule la perte
        conditionnelle au-del\`a de la VaR.
  \item \cs{Task<double[]> SimulatePortfolioReturns(IReadOnlyList<Position> positions,}\\
        \cs{MarketSnapshot snapshot, int horizon, CancellationToken ct)} -- retourne
        le vecteur brut des PnL simul\'es pour histogramme.
\end{itemize}

\begin{tcolorbox}[constraint]
\textbf{Parallel.For + ThreadLocal<Random> :} la boucle de simulation est CPU-bound.
Utilise \cs{Parallel.For} avec \cs{ThreadLocal<Random>} seed\'e par thread. D\'emontre
le speedup par rapport \`a une version s\'equentielle. Mesure et documente.\\[4pt]
\textbf{Interlocked sur les accumulateurs :} comme pour le pricer MC, l'accumulation
des PnL scenarios doit \^etre thread-safe. Utilise un tableau pr\'e-allou\'e index\'e par
scenario (pas d'accumulation partageable) ou \cs{Parallel.For} avec \cs{localFinally}.
\end{tcolorbox}

\begin{tcolorbox}[todo]
\textbf{Greeks analytiques multi-actifs (Vanna, Volga, Vanna crois\'ee).}
Pour les produits multi-actifs, les Greeks de second ordre (Vanna Equity/FX, Gamma
crois\'e) requi\`erent des m\'ethodes num\'eriques de second ordre ou des approches
pathwise / Malliavin plus avanc\'ees. Impl\'ementation pr\'evue avec l'introduction
des mod\`eles \`a volatilit\'e stochastique.
\end{tcolorbox}

\subsubsection{\texttt{VaRResult.cs} et \texttt{CVaRResult.cs}}

\textbf{Types :} deux \cs{record} immuables.

\cs{VaRResult} :
\begin{itemize}[leftmargin=1.5em]
  \item \cs{decimal VaR} -- perte maximale au niveau de confiance choisi.
  \item \cs{double ConfidenceLevel} -- ex: 0.99.
  \item \cs{int Horizon} -- en jours ouvrables.
  \item \cs{int ScenariosUsed}
  \item \cs{DateTimeOffset ComputedAt}
\end{itemize}

\cs{CVaRResult} :
\begin{itemize}[leftmargin=1.5em]
  \item Inclut toutes les propri\'et\'es de \cs{VaRResult}.
  \item \cs{decimal CVaR} -- perte conditionnelle esp\'er\'ee au-del\`a de la VaR.
  \item \cs{decimal[] TailLosses} -- distribution des pertes dans la queue.
\end{itemize}

% ============================================================
\subsection{PnL/ -- Moteur de calcul de PnL}

\subsubsection{\texttt{PnLEngine.cs}}

\textbf{R\^ole :} s'abonner aux mises \`a jour de march\'e et recalculer le PnL de
l'ensemble du portefeuille \`a chaque tick. Impl\'emente \cs{IObserver<MarketEnvironment>}
plut\^ot que \cs{IObserver<MarketSnapshot>}, afin que les Greeks agr\'eg\'es tiennent
compte des positions FX crois\'ees.

\textbf{M\'ethodes principales :}
\begin{itemize}[leftmargin=1.5em]
  \item \cs{void OnNext(MarketEnvironment env)} -- re\c{c}oit une nouvelle mise \`a jour
        d'environnement, recalcule le PnL de toutes les positions actives.
  \item \cs{void OnError(Exception error)} -- log l'erreur sans propagation.
  \item \cs{void OnCompleted()} -- flush final, sauvegarde en DB.
  \item \cs{Task<PnLSummary> GetCurrentPnLAsync(Guid positionId, CancellationToken ct)}
  \item \cs{Task<PortfolioPnL> GetPortfolioPnLAsync(CancellationToken ct)}
\end{itemize}

\textbf{Calcul du PnL :}
\[
  \text{PnL Daily} = Q \times (P_{\text{current}} - P_{\text{previous}})
\]
\[
  \text{PnL MTD} = Q \times (P_{\text{current}} - P_{\text{first of month}})
\]

\begin{tcolorbox}[constraint]
\textbf{Pattern Observer complet :} \cs{PnLEngine} impl\'emente \cs{IObserver<MarketEnvironment>}.
\cs{SimulatedMarketDataFeed} impl\'emente \cs{IObservable<MarketSnapshot>}. L'abonnement
via \cs{feed.Subscribe(pnlEngine)} retourne un \cs{IDisposable} stock\'e et appel\'e
dans le \cs{Dispose()} du \cs{PnLEngine}.\\[4pt]
\textbf{Thread-safety des events :} utiliser le pattern \cs{var h = PnLUpdated; h?.Invoke(...)}
pour \'eviter la race condition entre check-null et invocation.
\end{tcolorbox}

\subsubsection{\texttt{PnLSummary.cs} et \texttt{PortfolioPnL.cs}}

\cs{PnLSummary} (par position) :
\begin{itemize}[leftmargin=1.5em]
  \item \cs{Guid PositionId}, \cs{string Ticker}
  \item \cs{decimal DailyPnL}, \cs{decimal MTD}, \cs{decimal YTD}
  \item \cs{Greeks CurrentGreeks}, \cs{DateTimeOffset ComputedAt}
\end{itemize}

\cs{PortfolioPnL} (agr\'egat) :
\begin{itemize}[leftmargin=1.5em]
  \item \cs{decimal TotalDailyPnL}, \cs{decimal TotalMTD}, \cs{decimal TotalYTD}
  \item \cs{Greeks AggregatedGreeks}
  \item \cs{IReadOnlyList<PnLSummary> Positions}, \cs{DateTimeOffset AsOf}
\end{itemize}

% ============================================================
\subsection{Limits/ -- V\'erification des limites de risque}

\subsubsection{\texttt{LimitChecker.cs}}

\textbf{R\^ole :} s'abonner aux mises \`a jour PnL et d\'eclencher une alerte si une
limite est franchie. Impl\'emente \cs{IObserver<PortfolioPnL>}.

\textbf{Limites \`a v\'erifier :}

\begin{center}
\begin{tabularx}{\textwidth}{>{\ttfamily}l l X}
\toprule
\textbf{Limite} & \textbf{Seuil par d\'efaut} & \textbf{Comportement} \\
\midrule
DailyLossLimit   & $-50\,000$ EUR              & Alerte Critical si PnL Daily $<$ seuil \\
DeltaLimit       & $|\Delta_{\text{net}}| > 100$ & Alerte Warning \\
VegaLimit        & $|\text{Vega}_{\text{net}}| > 500$ & Alerte Warning \\
PositionSizeLimit & Notionnel $> 10$M EUR      & Alerte Info \\
VaRBreachLimit   & VaR $>$ 200\,000 EUR        & Alerte Critical (si VaR calcul\'ee) \\
\bottomrule
\end{tabularx}
\end{center}

\textbf{Comportements attendus :}
\begin{itemize}[leftmargin=1.5em]
  \item Une alerte ne doit \^etre d\'eclench\'ee qu'une seule fois par franchissement.
  \item Les seuils sont injectables via \cs{LimitConfiguration}.
  \item Chaque alerte est persist\'ee via \cs{IRepository<Alert>} inject\'e.
\end{itemize}

\begin{tcolorbox}[constraint]
\textbf{Test unitaire obligatoire :} l'alerte est d\'eclench\'ee exactement au
franchissement du seuil, pas avant, pas deux fois. Utilise \cs{Mock<IRepository<Alert>>}
(Moq) pour \'eviter la d\'ependance DB.
\end{tcolorbox}

% ============================================================
\subsection{StressTesting/ -- Sc\'enarios de choc}

\subsubsection{\texttt{StressTestEngine.cs}}

\textbf{R\^ole :} \'evaluer l'impact de chocs de march\'e sur le portefeuille.
Chaque sc\'enario tourne en parall\`ele via \cs{Task.WhenAll}.

\textbf{Sc\'enarios requis :}

\begin{center}
\begin{tabularx}{\textwidth}{l l X}
\toprule
\textbf{Sc\'enario}    & \textbf{Param\`etres}                        & \textbf{Description} \\
\midrule
SpotShock             & $-5\%$, $-10\%$, $-20\%$, $-30\%$           & Baisse du spot \\
VolShock              & $+5$~pts, $+10$~pts, $+20$~pts               & Hausse de la volatilit\'e \\
CombinedShock         & $-10\%$ spot $+$ $+10$~pts vol               & Sc\'enario combin\'e \\
TailRisk              & $-40\%$ spot $+$ $+30$~pts vol               & Sc\'enario de queue \\
CorrelationShock      & $\Delta\rho = +0.20$ sur toutes les paires   & Choc de corr\'elation \\
FXShock               & $\pm 10\%$ sur EUR/USD                       & Choc de change \\
\bottomrule
\end{tabularx}
\end{center}

\textbf{M\'ethodes principales :}
\begin{itemize}[leftmargin=1.5em]
  \item \cs{Task<IReadOnlyList<StressTestResult>> RunAllScenariosAsync(}\\
        \cs{IReadOnlyList<Position> positions, IProgress<StressTestProgress> progress,}\\
        \cs{CancellationToken ct)}
  \item \cs{Task<StressTestResult> RunScenarioAsync(StressScenario scenario,}\\
        \cs{IReadOnlyList<Position> positions, CancellationToken ct)}
  \item \cs{StressTestMatrix BuildMatrix(IReadOnlyList<StressTestResult> results)}\\
        via LINQ uniquement (\cs{GroupBy}, \cs{SelectMany}).
\end{itemize}

\begin{tcolorbox}[constraint]
\textbf{Task.WhenAll obligatoire :} chaque sc\'enario est une \cs{Task} ind\'ependante.\\[4pt]
\textbf{IProgress<T> :} la progression rapport\'ee via \cs{IProgress<StressTestProgress>}.\\[4pt]
\textbf{LINQ avanc\'e :} agr\'egation en matrice via \cs{GroupBy}/\cs{SelectMany}.
Interdiction de boucles \cs{foreach} pour les transformations de donn\'ees.
\end{tcolorbox}

\subsubsection{\texttt{StressTestResult.cs} et \texttt{StressTestProgress.cs}}

\cs{StressTestResult} :
\begin{itemize}[leftmargin=1.5em]
  \item \cs{StressScenario Scenario}, \cs{Guid PositionId}, \cs{string Ticker}
  \item \cs{decimal PnLImpact}, \cs{Greeks GreeksUnderShock}, \cs{DateTimeOffset ComputedAt}
\end{itemize}

\cs{StressTestProgress} :
\begin{itemize}[leftmargin=1.5em]
  \item \cs{int CompletedScenarios}, \cs{int TotalScenarios}, \cs{string CurrentScenarioName}
  \item \cs{double PercentComplete} -- propri\'et\'e calcul\'ee.
\end{itemize}

% ============================================================
\subsection{Feeds/ -- Simulation du flux de march\'e temps r\'eel}

\subsubsection{\texttt{SimulatedMarketDataFeed.cs}}

\textbf{R\^ole :} publier un \cs{MarketSnapshot} simul\'e toutes les secondes.
Impl\'emente \cs{IObservable<MarketSnapshot>} et \cs{IMarketDataFeed}.

\textbf{Comportements attendus :}
\begin{itemize}[leftmargin=1.5em]
  \item Prix spot simul\'e : $S_{t+1} = S_t \cdot (1 + \sigma\sqrt{dt} \cdot Z)$,
        $Z \sim \mathcal{N}(0,1)$.
  \item \cs{IDisposable Subscribe(IObserver<MarketSnapshot> observer)} -- d\'esabonnement
        via \cs{Dispose()}.
  \item \cs{Task StartAsync(CancellationToken ct)} -- boucle de publication
        (\cs{Task.Delay(1000, ct)}).
\end{itemize}

\begin{tcolorbox}[constraint]
\textbf{Memory leak et Dispose :} d\'emontre le leak m\'emoire si l'abonnement n'est
pas d\'esabonné (\cs{Dispose} sur le token de souscription), puis montre la correction.\\[4pt]
\textbf{Event thread-safe :} pattern \cs{var snapshot = MarketUpdated; snapshot?.Invoke(...)}
pour la notification des abonn\'es delegates.
\end{tcolorbox}

\begin{tcolorbox}[deliverable]
Risk v2.0 est termin\'e quand :
\begin{itemize}[leftmargin=1.5em]
  \item CVaR $\ge$ VaR pour tout niveau de confiance et tout portefeuille test.
  \item Un portefeuille long A / short B avec $\rho_{AB} = 0.9$ a une VaR
        strictement inf\'erieure \`a $\text{VaR}(A) + \text{VaR}(B)$
        (b\'en\'efice de diversification).
  \item Un portefeuille long A / short B avec $\rho_{AB} = -1$ a une VaR proche de 0
        (couverture parfaite).
  \item \cs{PnLEngine} ne provoque aucun memory leak (test de d\'esabonnement Moq).
  \item Les sc\'enarios de choc de corr\'elation sont ex\'ecut\'es en parall\`ele
        via \cs{Task.WhenAll} et rapportent la progression via \cs{IProgress<T>}.
\end{itemize}
\end{tcolorbox}

\newpage
\section{TradingDashboard.Data}
% ------------------------------------------------------------

\begin{tcolorbox}[note]
Ce projet est la couche de persistance de toute l'application. Il expose les données
via des repositories et ne contient aucune logique métier. Toute règle financière
appartient à \texttt{Core} ou \texttt{Risk}, pas ici.
\end{tcolorbox}

\subsection{TradingDbContext.cs -- Contexte Entity Framework Core}

\textbf{Rôle} : point d'entrée unique vers la base de données SQLite. Hérite de
\cs{DbContext} (EF Core 8).

\textbf{Propriétés \cs{DbSet<T>} attendues :}
\begin{itemize}
  \item \cs{DbSet<Position> Positions}
  \item \cs{DbSet<Instrument> Instruments}
  \item \cs{DbSet<Trade> Trades}
  \item \cs{DbSet<MarketData> MarketDataHistory}
  \item \cs{DbSet<PnLSnapshot> PnLHistory}
  \item \cs{DbSet<Alert> Alerts}
\end{itemize}

\textbf{Comportements attendus :}
\begin{itemize}
  \item Surcharger \cs{OnConfiguring} pour pointer vers \texttt{trading.db} (SQLite).
  \item Surcharger \cs{OnModelCreating} pour configurer les relations et les
        conversions de type (voir contraintes).
  \item Configurer la relation 1-N entre \cs{Instrument} et \cs{Position} :
        une position est obligatoirement liée à un instrument existant (clé étrangère
        non nullable).
  \item Configurer la relation 1-N entre \cs{Position} et \cs{PnLSnapshot} :
        la suppression d'une position doit déclencher la suppression en cascade de
        ses snapshots (\cs{DeleteBehavior.Cascade}).
  \item Configurer la relation 1-N entre \cs{Position} et \cs{Alert} avec
        \cs{DeleteBehavior.SetNull} (une alerte orpheline doit être conservée en
        audit trail).
  \item Stocker les Greeks (struct) en colonne JSON via un \emph{value converter}
        EF Core. La colonne s'appelle \cs{GreeksJson} de type \cs{TEXT}.
\end{itemize}

\begin{tcolorbox}[constraint]
\textbf{IDisposable et cycle de vie :} \cs{TradingDbContext} implémente déjà
\cs{IDisposable} via \cs{DbContext}. Partout où tu l'instancies manuellement
(hors injection de dépendances), utilise un bloc \cs{using}. Comprends pourquoi
l'oubli du \cs{Dispose} laisse la connexion SQLite ouverte et bloque les accès
concurrents. \\[4pt]
\textbf{Logging SQL :} Active \cs{EnableSensitiveDataLogging()} et
\cs{LogTo(Console.WriteLine)} uniquement en mode Debug. Tu devras montrer le SQL
généré par EF lors des entretiens -- c'est non négociable.
\end{tcolorbox}

% ------------------------------------------------------------
\subsection{Migrations/ -- Migrations EF Core}

\textbf{Règle absolue :} aucune modification manuelle de \texttt{trading.db}.
Toute évolution du schéma passe par une migration nommée explicitement.

\textbf{Migrations à créer dans cet ordre :}
\begin{enumerate}
  \item \texttt{InitialCreate} -- tables de base : Instruments, Positions, Trades.
  \item \texttt{AddMarketDataHistory} -- table MarketData avec index sur
        \cs{(InstrumentId, Timestamp)}.
  \item \texttt{AddPnLAndAlerts} -- tables PnLSnapshot et Alerts, colonne
        \cs{GreeksJson}.
  \item \texttt{AddAlertAcknowledgedColumn} -- ajout de \cs{IsAcknowledged} et
        \cs{AcknowledgedAt} sur la table Alerts.
\end{enumerate}

\begin{tcolorbox}[note]
Chaque migration génère deux méthodes : \cs{Up()} et \cs{Down()}. Tu dois être
capable d'expliquer ce que fait \cs{Down()} et dans quel cas elle est invoquée.
C'est une question d'entretien classique sur EF Core.
\end{tcolorbox}

% ------------------------------------------------------------
\subsection{Repositories/ -- Pattern Repository}

\subsubsection{GenericRepository.cs}

Implémente l'interface \cs{IRepository<T>} définie dans \texttt{Core}.
\cs{T} est contraint par \cs{where T : class}.

\textbf{Méthodes à implémenter :}
\begin{itemize}
  \item \cs{Task<T?> GetByIdAsync(Guid id, CancellationToken ct = default)}
  \item \cs{Task<IReadOnlyList<T>> GetAllAsync(CancellationToken ct = default)}
  \item \cs{Task AddAsync(T entity, CancellationToken ct = default)}
  \item \cs{Task UpdateAsync(T entity, CancellationToken ct = default)}
  \item \cs{Task DeleteAsync(Guid id, CancellationToken ct = default)}
  \item \cs{Task SaveChangesAsync(CancellationToken ct = default)}
\end{itemize}

\begin{tcolorbox}[constraint]
\textbf{IQueryable obligatoire :} les méthodes qui filtrent des données doivent
retourner \cs{IQueryable<T>} et non \cs{IEnumerable<T>}. La différence est
fondamentale : \cs{IQueryable} compose la requête et la traduit en SQL ;
\cs{IEnumerable} charge tout en mémoire puis filtre côté application (N+1
problem). Tu devras montrer les deux comportements avec le logging SQL activé.
\end{tcolorbox}

\subsubsection{TradeRepository.cs}

Hérite de \cs{GenericRepository<Trade>} et ajoute des requêtes métier.

\textbf{Méthodes métier spécifiques :}
\begin{itemize}
  \item \cs{Task<IReadOnlyList<Trade>> GetTradesByDeskAsync(string desk, CancellationToken ct)}\\
        Filtre les trades par desk. Doit utiliser \cs{IQueryable} et LINQ.
  \item \cs{Task<IReadOnlyList<Trade>> GetTradesForInstrumentAsync(Guid instrumentId, CancellationToken ct)}\\
        Inclure les données de l'instrument associé (\emph{eager loading} via
        \cs{Include}).
  \item \cs{Task<decimal> GetTotalNotionalByDeskAsync(string desk, CancellationToken ct)}\\
        Agrégation SQL : utilise \cs{SumAsync} de EF Core, pas de Sum LINQ
        sur une liste en mémoire.
\end{itemize}

\subsubsection{PositionRepository.cs}

\textbf{Méthodes métier spécifiques :}
\begin{itemize}
  \item \cs{Task<IReadOnlyList<Position>> GetOpenPositionsAsync(CancellationToken ct)}\\
        Filtre les positions dont le statut est \cs{Active}. Inclut l'instrument
        et les 30 derniers snapshots PnL (eager loading limité).
  \item \cs{Task<IReadOnlyList<Position>> GetPositionsByUnderlyingAsync(string ticker, CancellationToken ct)}\\
        Joint sur la table Instruments pour filtrer par sous-jacent.
  \item \cs{Task<Greeks> GetAggregatedGreeksAsync(CancellationToken ct)}\\
        Somme les Greeks de toutes les positions actives. L'agrégation doit se
        faire en mémoire après un seul appel DB (les Greeks sont en JSON, EF ne
        peut pas les agréger en SQL).
\end{itemize}

\textbf{Démonstration obligatoire -- Eager vs Lazy Loading :}

Tu dois implémenter et documenter les deux comportements sur
\cs{GetOpenPositionsAsync} :
\begin{itemize}
  \item Version eager : \cs{.Include(p => p.Instrument).Include(p => p.PnLHistory)} --
        1 requête SQL avec JOIN.
  \item Version lazy : accès à \cs{position.Instrument} sans \cs{Include} --
        déclenche N requêtes supplémentaires (N+1 problem).
\end{itemize}
Le logging SQL doit être activé pour visualiser la différence.

\subsubsection{AlertRepository.cs}

\textbf{Méthodes métier spécifiques :}
\begin{itemize}
  \item \cs{Task<IReadOnlyList<Alert>> GetActiveAlertsAsync(CancellationToken ct)}\\
        Filtre \cs{IsAcknowledged == false}, tri par \cs{Level} décroissant.
  \item \cs{Task AcknowledgeAlertAsync(Guid alertId, CancellationToken ct)}\\
        Met à jour \cs{IsAcknowledged = true} et \cs{AcknowledgedAt = DateTime.UtcNow}.
        Utilise une mise à jour ciblée (\cs{ExecuteUpdateAsync}) sans charger
        l'entité entière.
  \item \cs{Task<IReadOnlyList<Alert>> GetAlertsByLevelAsync(AlertLevel level, CancellationToken ct)}
\end{itemize}

\begin{tcolorbox}[constraint]
\textbf{LINQ obligatoire sur toutes les requêtes métier.} Zéro SQL raw sauf
justification documentée en commentaire XML. Si tu utilises \cs{FromSqlRaw},
explique pourquoi LINQ ne permettait pas d'exprimer cette requête. \\[4pt]
\textbf{async/await partout :} toutes les méthodes qui touchent la DB sont
\cs{async Task<T>}. Interdiction de \cs{.Result} ou \cs{.Wait()}. Utilise
\cs{ConfigureAwait(false)} dans les méthodes des repositories (couche librairie,
pas UI).
\end{tcolorbox}

% ------------------------------------------------------------
\subsection{Schéma de la base de données}

\begin{center}
\begin{tabularx}{\textwidth}{>{\ttfamily}l >{\ttfamily}l X}
\toprule
\textbf{Table} & \textbf{Colonne} & \textbf{Contrainte / Note} \\
\midrule
Instruments & Id (GUID) & Clé primaire, généré automatiquement \\
            & Ticker & NOT NULL, longueur max 20 \\
            & ProductType & Enum stocké en string (HasConversion) \\
            & Strike & REAL nullable \\
            & Underlying & NOT NULL \\
            & Currency & Enum stocké en string \\
            & Maturity & TEXT (DateOnly sérialisé) \\
\midrule
Positions & Id (GUID) & Clé primaire \\
          & InstrumentId & Clé étrangère NOT NULL vers Instruments \\
          & NetQuantity & INTEGER \\
          & Status & Enum (Active / Closed) en string \\
          & OpenDate & TEXT (DateTimeOffset) \\
\midrule
Trades & TradeId (GUID) & Clé primaire \\
       & InstrumentId & Clé étrangère NOT NULL \\
       & Notional & REAL NOT NULL \\
       & Quantity & INTEGER NOT NULL \\
       & TradePrice & REAL NOT NULL \\
       & TradeDate & TEXT (DateTimeOffset ISO8601) \\
       & Desk & NOT NULL, longueur max 50 \\
       & Trader & NOT NULL \\
\midrule
MarketData & Id (GUID) & Clé primaire \\
           & InstrumentId & Clé étrangère \\
           & SpotPrice & REAL NOT NULL \\
           & ImpliedVolatility & REAL \\
           & RiskFreeRate & REAL \\
           & Timestamp & TEXT (DateTimeOffset) \\
           & Source & NOT NULL \\
\midrule
PnLHistory & Id (GUID) & Clé primaire \\
           & PositionId & Clé étrangère, cascade delete \\
           & DailyPnL & REAL \\
           & MTD & REAL \\
           & YTD & REAL \\
           & GreeksJson & TEXT (serialisé JSON) \\
           & Date & TEXT (DateOnly) \\
\midrule
Alerts & Id (GUID) & Clé primaire \\
       & Level & Enum en string \\
       & Message & TEXT NOT NULL \\
       & PositionId & Clé étrangère nullable (SetNull) \\
       & TriggeredAt & TEXT (DateTimeOffset) \\
       & IsAcknowledged & INTEGER (bool) \\
       & AcknowledgedAt & TEXT nullable \\
\bottomrule
\end{tabularx}
\end{center}


\newpage
% ------------------------------------------------------------
\section{Tests unitaires associés}
% ------------------------------------------------------------

\subsection{Tests pour .Data}

\begin{itemize}
  \item \textbf{TradeRepositoryTests} : CRUD complet sur SQLite in-memory
        (\cs{UseInMemoryDatabase}). Vérification que \cs{GetTradesByDeskAsync}
        filtre correctement.
  \item \textbf{PositionRepositoryTests} : vérifier le cascade delete sur
        PnLSnapshot. Vérifier que \cs{GetOpenPositionsAsync} n'inclut pas les
        positions closes.
  \item \textbf{AlertRepositoryTests} : vérifier que \cs{AcknowledgeAlertAsync}
        ne déclenche pas de chargement complet de l'entité (utiliser le logging EF
        pour compter les requêtes SQL).
\end{itemize}

\subsection{Tests pour .Risk}

\begin{itemize}
  \item \textbf{MonteCarloEngineTests} : prix MC dans $\pm 2\%$ du prix BS pour
        100\,000 chemins. Test de convergence : vérifier que l'erreur standard
        diminue en $1/\sqrt{N}$.
  \item \textbf{PnLEngineTests} : simuler un mouvement spot de $+1\%$ et vérifier
        que le PnL Daily d'une position longue call est positif et cohérent avec
        le Delta.
  \item \textbf{LimitCheckerTests} : vérifier le franchissement exact du seuil,
        l'absence de doublon d'alerte, la persistance via \cs{Mock<IRepository<Alert>>}.
  \item \textbf{StressTestEngineTests} : vérifier que tous les scénarios tournent
        en parallèle (mesurer que le temps total est inférieur à la somme des
        temps individuels).
\end{itemize}

\newpage

% ------------------------------------------------------------
\section{TradingDashboard.Reports}
% ------------------------------------------------------------

\begin{tcolorbox}[note]
Ce projet produit les exports de risque journaliers en Excel (EPPlus) et PDF
(QuestPDF). Il consomme les données via les interfaces \cs{IRepository<T>}
et le \cs{PnLEngine}. Aucune logique de calcul ici -- uniquement mise en forme
et export.
\end{tcolorbox}

% ------------------------------------------------------------
\subsection{ExcelReporter.cs -- Export Excel}

\textbf{Rôle} : générer un classeur Excel multi-onglets représentant
l'état complet du portefeuille. Utilise \textbf{EPPlus} (licence
NonCommercial gratuite pour usage éducatif).

\textbf{Implémente} \cs{IDisposable} -- le \cs{ExcelPackage} EPPlus doit
être libéré dans \cs{Dispose()}.

\textbf{Onglets à générer :}

\textbf{Onglet 1 -- Positions :}
\begin{itemize}
  \item Colonnes : Ticker, ProductType, NetQuantity, Strike, Maturity,
        DaysToMaturity, Currency, Status.
  \item Mise en forme : en-têtes en gras avec fond vert foncé (\texttt{\#008264}),
        lignes alternées gris clair, filtre automatique activé.
  \item Tri : par DaysToMaturity croissant (positions proche échéance en premier).
  \item Formule Excel générée programmatiquement : colonne
        ``Valeur Notionnelle EUR'' = \texttt{=NetQuantity * Strike}.
\end{itemize}

\textbf{Onglet 2 -- PnL History :}
\begin{itemize}
  \item Colonnes : Date, Ticker, DailyPnL, MTD, YTD, Delta, Vega, Theta.
  \item Mise en forme conditionnelle : DailyPnL en rouge si négatif, en
        vert si positif (via \cs{ExcelConditionalFormattingIconDataBar}).
  \item Ligne de total en bas : somme des colonnes numériques via formule
        \texttt{SUM} générée dynamiquement.
\end{itemize}

\textbf{Onglet 3 -- Greeks :}
\begin{itemize}
  \item Tableau des Greeks par position et Greeks agrégés du portefeuille.
  \item Graphique en barres : Delta net par sous-jacent, généré
        programmatiquement via \cs{ExcelChart}.
\end{itemize}

\textbf{Méthodes :}
\begin{itemize}
  \item \cs{Task<byte[]> GenerateReportAsync(PortfolioPnL portfolio, CancellationToken ct = default)}\\
        Retourne le fichier Excel en mémoire (byte array). Pas d'écriture
        disque directe -- l'appelant décide où sauvegarder.
  \item \cs{Task SaveToFileAsync(byte[] content, string path, CancellationToken ct = default)}
\end{itemize}

\begin{tcolorbox}[constraint]
\textbf{IDisposable et using :} \cs{ExcelReporter} implémente
\cs{IDisposable}. Partout où il est instancié manuellement, il doit
être dans un bloc \cs{using}. L'\cs{ExcelPackage} interne doit être
libéré dans \cs{Dispose()} même si \cs{GenerateReportAsync}
n'a pas été appelé. \\[4pt]
\textbf{LINQ pour la préparation des données :} toutes les
transformations (tri, filtrage, projection vers les colonnes Excel) doivent
être en LINQ. Pas de boucle \cs{foreach} pour construire les tableaux de
données.
\end{tcolorbox}

% ------------------------------------------------------------
\subsection{PdfReporter.cs -- Export PDF}

\textbf{Rôle} : générer un rapport de risque journalier mis en page
via \textbf{QuestPDF} (open-source, licence MIT pour usage éducatif).

\textbf{Implémente} \cs{IDisposable}.

\textbf{Structure du rapport PDF :}
\begin{enumerate}
  \item \textbf{Page de garde} : titre ``Rapport de
        Risque Journalier'', date de génération, période couverte.
  \item \textbf{Résumé exécutif} (1 page) : PnL Daily total,
        PnL MTD, PnL YTD, Greeks agrégés du portefeuille, nombre
        d'alertes actives. Mise en forme tabulée avec code couleur
        vert/rouge selon signe du PnL.
  \item \textbf{Détail par position} (N pages) : une section par position
        active avec : identité de l'instrument, Greeks individuels, historique
        PnL sur 30 jours en tableau.
  \item \textbf{Alertes actives} : tableau des alertes non acquittées trié
        par sévérité décroissante.
  \item \textbf{Stress Test Summary} : matrice des impacts PnL par scénario
        (si les résultats sont disponibles).
\end{enumerate}

\textbf{Méthodes :}
\begin{itemize}
  \item \cs{Task<byte[]> GenerateReportAsync(PortfolioPnL portfolio, IReadOnlyList<Alert> alerts, CancellationToken ct = default)}
  \item \cs{Task SaveToFileAsync(byte[] content, string path, CancellationToken ct = default)}
\end{itemize}

\begin{tcolorbox}[note]
QuestPDF utilise un modèle de composition déclaratif (similaire à
Jetpack Compose ou SwiftUI) : tu décris le document en C\# pur via des
méthodes chaînées (\cs{.Row()}, \cs{.Column()},
\cs{.Table()}). C'est un exercice d'utilisation d'une API fluent -- pattern
Builder appliqué à la génération de documents.
\end{tcolorbox}

% ------------------------------------------------------------
\subsection{ReportScheduler.cs -- Planification}

\textbf{Rôle} : déclencher la génération automatique des rapports
à heure fixe (17h00 heure de Paris, fin de journée de trading).

\textbf{Implémentation :}
\begin{itemize}
  \item Utilise \cs{PeriodicTimer} pour vérifier toutes les minutes si
        l'heure de génération est atteinte.
  \item Génère Excel et PDF en parallèle via
        \cs{Task.WhenAll(excelTask, pdfTask)}.
  \item Sauvegarde les fichiers dans un dossier \cs{./reports/YYYY-MM-DD/}.
  \item Log la durée de génération et la taille des fichiers produits.
\end{itemize}

\begin{tcolorbox}[constraint]
\textbf{CancellationToken propagé :} \cs{ReportScheduler} accepte un
\cs{CancellationToken} dans sa méthode \cs{StartAsync}. Quand le token
est annulé (arrêt de l'application), la boucle se termine proprement sans
exception non gérée. Utilise \cs{PeriodicTimer.WaitForNextTickAsync(ct)}
qui lance une \cs{OperationCanceledException} que tu catches pour sortir
proprement.
\end{tcolorbox}

\newpage

% ------------------------------------------------------------
\section{TradingDashboard.Dashboard}
% ------------------------------------------------------------

\begin{tcolorbox}[note]
C'est le point d'entrée de l'application (\cs{Program.cs}). Il orchestre
tous les modules, configure l'injection de dépendances, et affiche l'état
du portefeuille en temps réel dans le terminal. C'est une \textbf{Console App}
\cs{.NET 8}.
\end{tcolorbox}

% ------------------------------------------------------------
\subsection{Program.cs -- Point d'entrée et DI}

\textbf{Rôle} : configurer le conteneur d'injection de dépendances
(\cs{ServiceCollection}), démarrer les services et lancer la boucle
d'affichage.

\textbf{Enregistrements DI attendus :}
\begin{itemize}
  \item \cs{AddDbContext<TradingDbContext>} -- cycle de vie \cs{Scoped}.
  \item \cs{AddScoped<IRepository<Trade>, TradeRepository>}
  \item \cs{AddScoped<IRepository<Position>, PositionRepository>}
  \item \cs{AddScoped<IRepository<Alert>, AlertRepository>}
  \item \cs{AddSingleton<IMarketDataFeed, MarketDataService>}
  \item \cs{AddSingleton<IPricer<VanillaOption, PricingResult>, BlackScholesPricer>}
  \item \cs{AddSingleton<PnLEngine>}
  \item \cs{AddSingleton<LimitChecker>}
  \item \cs{AddSingleton<MarketDataCache>}
  \item \cs{AddHttpClient<AlphaVantageClient>} avec politique Polly attachée.
\end{itemize}

\textbf{Séquence de démarrage :}
\begin{enumerate}
  \item Appliquer les migrations EF Core en attente (\cs{context.Database.MigrateAsync()}).
  \item Injecter des données de test si la base est vide (seed data :
        5 instruments, 10 trades répartis sur 2 desks).
  \item Démarrer le \cs{SimulatedMarketDataFeed} (ou l'API réelle selon
        la configuration).
  \item Abonner \cs{PnLEngine} au feed.
  \item Abonner \cs{LimitChecker} au \cs{PnLEngine}.
  \item Lancer la boucle d'affichage console.
\end{enumerate}

\begin{tcolorbox}[constraint]
\textbf{Top-level statements :} \cs{Program.cs} doit utiliser les top-level
statements (C\# 9+). Pas de classe \cs{Program} explicite ni de méthode
\cs{Main}. Sache expliquer ce que le compilateur génère derrière
cette syntaxe. \\[4pt]
\textbf{CancellationToken global :} créer un \cs{CancellationTokenSource}
principal relié à \cs{Console.CancelKeyPress} (Ctrl+C). Ce token est
propagé à tous les services au démarrage. L'arrêt doit être
propre~: flush du PnL, sauvegarde en base, fermeture du feed.
\end{tcolorbox}

% ------------------------------------------------------------
\subsection{ConsoleRenderer.cs -- Affichage temps réel}

\textbf{Rôle} : rafraîchir l'affichage console toutes les secondes avec
l'état courant du portefeuille. Utilise \cs{Console.SetCursorPosition} pour
éviter le scroll (affichage en place fixe).

\textbf{Sections à afficher :}

\textbf{Section 1 -- En-tête :}
\begin{itemize}
  \item Heure courante, version de l'application, statut du feed
        (Connected / Simulated / Error).
  \item Barre de séparation horizontale.
\end{itemize}

\textbf{Section 2 -- PnL Summary :}

\begin{center}
\begin{tabular}{lrrr}
\toprule
\textbf{Desk} & \textbf{Daily PnL} & \textbf{MTD} & \textbf{YTD} \\
\midrule
DERIV-EU-01 & +12 450 EUR & +87 200 EUR & +342 000 EUR \\
Portfolio Total & +12 450 EUR & +87 200 EUR & +342 000 EUR \\
\bottomrule
\end{tabular}
\end{center}

Couleur : vert (\cs{ConsoleColor.Green}) si positif, rouge
(\cs{ConsoleColor.Red}) si négatif. Utilise \cs{Console.ForegroundColor}
et restaure-le après chaque ligne.

\textbf{Section 3 -- Greeks agrégés :}
\begin{itemize}
  \item Delta net, Gamma net, Vega net, Theta du portefeuille.
  \item Alerte visuelle (\cs{[!]}) si un Greek dépasse sa limite.
\end{itemize}

\textbf{Section 4 -- Blotter de positions :}
\begin{itemize}
  \item Tableau des positions actives : Ticker, Type, Qty, Strike, DaysToMaturity,
        Prix Théorique, Delta, PnL Daily.
  \item Tri par DaysToMaturity croissant via LINQ (\cs{.OrderBy}).
  \item Positions proches de l'échéance ($<$ 5 jours) affichées en
        jaune (\cs{ConsoleColor.Yellow}).
\end{itemize}

\textbf{Section 5 -- Queue d'alertes :}
\begin{itemize}
  \item Les 5 dernières alertes actives triées par sévérité
        décroissante via \cs{PriorityQueue<Alert, AlertLevel>}.
  \item Alertes \cs{Critical} et \cs{Fatal} en rouge clignotant (simulé
        via alternance d'affichage).
\end{itemize}

\begin{tcolorbox}[constraint]
\textbf{SortedDictionary et PriorityQueue :} l'historique PnL des 60 derniers
ticks doit être stocké dans un
\cs{SortedDictionary<DateTimeOffset, PortfolioPnL>} (accès trié
chronologiquement, O(log n) en insertion). La queue d'alertes utilise
\cs{PriorityQueue<Alert, AlertLevel>} (min-heap). Pour avoir les alertes
les plus critiques en premier, inverse la comparaison via un
\cs{IComparer<AlertLevel>} personnalisé. \\[4pt]
\textbf{LinkedList comme circular buffer :} les snapshots de marché en
mémoire (1000 derniers) sont stockés dans une
\cs{LinkedList<MarketSnapshot>}. Quand la taille dépasse 1000,
supprime le premier élément (\cs{RemoveFirst()}) avant d'ajouter le
nouveau (\cs{AddLast()}). Sache expliquer pourquoi \cs{LinkedList} est
préférable à \cs{List<T>} pour ce pattern (O(1) en ajout/suppression
aux extrémités).
\end{tcolorbox}

% ------------------------------------------------------------
\subsection{Configuration/ -- Paramètres de l'application}

\textbf{Fichier} \cs{appsettings.json} à créer à la racine du projet
Dashboard. Doit contenir :

\begin{itemize}
  \item \cs{AlphaVantage:ApiKey} -- clé API (ne jamais committer en clair,
        utiliser \cs{dotnet user-secrets} pour le développement local).
  \item \cs{MarketData:CacheExpirySeconds} -- durée de vie du cache (30 par
        défaut).
  \item \cs{MarketData:UseSimulation} -- bool. Si true, utilise
        \cs{SimulatedMarketDataFeed} au lieu de l'API réelle.
  \item \cs{Risk:DailyLossLimit} -- seuil en EUR (-50000 par défaut).
  \item \cs{Risk:DeltaLimit}, \cs{Risk:VegaLimit}.
  \item \cs{Dashboard:RefreshIntervalMs} -- intervalle de rafraîchissement
        console (1000 par défaut).
\end{itemize}

\newpage

% ------------------------------------------------------------
\section{TradingDashboard.Tests}
% ------------------------------------------------------------

\begin{tcolorbox}[note]
Projet transversal xUnit couvrant tous les modules. La couverture minimale
requise est 70\% sur les projets \texttt{Pricing} et \texttt{Risk}. Ce projet
référence tous les autres projets (sauf \texttt{Dashboard}).
\end{tcolorbox}

% ------------------------------------------------------------
\subsection{Pricing/ -- Tests du moteur de pricing}

\subsubsection{BlackScholesTests.cs}

\textbf{Cas de test obligatoires :}
\begin{itemize}
  \item \textbf{Prix call ATM} : $S=100$, $K=100$, $\sigma=20\%$, $r=5\%$,
        $T=1$ an. Résultat attendu : $\approx 10.45$. Tolérance : $\pm 0.01$.
  \item \textbf{Put-Call Parity} : vérifier que
        $C - P = S \cdot e^{-qT} - K \cdot e^{-rT}$ à $10^{-6}$ près
        pour tout jeu de paramètres.
  \item \textbf{Delta call deep ITM} : $S=150$, $K=100$, $T=0.1$. Delta doit
        être $> 0.99$.
  \item \textbf{Delta call deep OTM} : $S=50$, $K=100$, $T=0.1$. Delta doit
        être $< 0.01$.
  \item \textbf{Gamma toujours positif} : vérifier sur 20 jeux de paramètres
        aléatoires (propriété-based testing avec une boucle paramétrique
        xUnit via \cs{[Theory][InlineData]}).
  \item \textbf{Theta négatif} : pour une position longue, Theta doit être
        négatif (décroissance de la valeur temporelle).
  \item \textbf{Vega positif} : Vega est toujours positif pour un call ou un put
        long (une hausse de vol augmente la valeur de l'option).
\end{itemize}

\subsubsection{MonteCarloValidationTests.cs}

\textbf{Cas de test obligatoires :}
\begin{itemize}
  \item \textbf{Convergence vers Black-Scholes} : prix MC à $\pm 2\%$ du
        prix BS pour 100\,000 chemins. Paramètres identiques aux tests BS.
  \item \textbf{Décroissance de l'erreur standard} : vérifier que
        $\text{StdError}(N=100\,000) < \text{StdError}(N=10\,000)$.
  \item \textbf{Speedup mesurable} : le mode parallèle doit être au moins
        1.5x plus rapide que le mode séquentiel pour 500\,000 chemins sur une
        machine multi-coeurs. Ce test peut être marqué
        \cs{[Fact(Skip = "Performance test -- run manually")]} pour ne pas
        alourdir la CI.
\end{itemize}

% ------------------------------------------------------------
\subsection{Data/ -- Tests de la couche de persistance}

\subsubsection{TradeRepositoryTests.cs}

\textbf{Setup} : utiliser \cs{UseInMemoryDatabase("TestDb\_" + Guid.NewGuid())}
pour isoler chaque test.

\textbf{Cas de test obligatoires :}
\begin{itemize}
  \item CRUD complet : Add, GetById, Update, Delete sur un Trade.
  \item \cs{GetTradesByDeskAsync} : insérer 10 trades sur 2 desks, vérifier
        que le filtre retourne exactement les trades du bon desk.
  \item \cs{GetTotalNotionalByDeskAsync} : vérifier l'agrégation
        arithmétique sur des données connues.
  \item Trade inexistant : \cs{GetByIdAsync} avec un Guid aléatoire doit
        retourner \cs{null} sans exception.
\end{itemize}

\subsubsection{PositionRepositoryTests.cs}

\textbf{Cas de test obligatoires :}
\begin{itemize}
  \item \cs{GetOpenPositionsAsync} ne retourne pas les positions avec statut
        \cs{Closed}.
  \item Cascade delete : supprimer une Position doit supprimer ses
        \cs{PnLSnapshot} associés.
  \item \cs{GetAggregatedGreeksAsync} : vérifier la somme sur 3 positions
        avec des Greeks connus.
\end{itemize}

\subsubsection{AlertRepositoryTests.cs}

\textbf{Cas de test obligatoires :}
\begin{itemize}
  \item \cs{AcknowledgeAlertAsync} : vérifier que \cs{IsAcknowledged}
        passe à true et que \cs{AcknowledgedAt} est renseigné.
  \item \cs{GetActiveAlertsAsync} : une alerte acquittée ne doit pas
        apparaître dans les actives.
  \item Tri par sévérité : vérifier que la première alerte
        retournée est bien la plus critique.
\end{itemize}

% ------------------------------------------------------------
\subsection{Risk/ -- Tests du moteur de risque}

\subsubsection{PnLEngineTests.cs}

\textbf{Setup} : mocker \cs{IRepository<Position>} et
\cs{IPricer<VanillaOption, PricingResult>} via Moq.

\textbf{Cas de test obligatoires :}
\begin{itemize}
  \item \textbf{Mouvement spot positif} : simuler un snapshot avec spot
        $+1\%$. Le PnL Daily d'une position longue call doit être positif.
        Cohérence avec Delta : $\Delta PnL \approx \Delta \times \Delta S \times Q$.
  \item \textbf{Mouvement spot négatif} : position longue call, spot $-1\%$,
        PnL Daily négatif.
  \item \textbf{Position short call} : Quantity $< 0$, spot $+1\%$, PnL Daily
        négatif (position courte perd quand le marché monte).
  \item \textbf{OnError ne propage pas} : appeler \cs{OnError} ne doit pas
        faire planter la boucle d'observation.
\end{itemize}

\subsubsection{LimitCheckerTests.cs}

\textbf{Setup} : \cs{Mock<IRepository<Alert>>} pour capturer les alertes
persistees.

\textbf{Cas de test obligatoires :}
\begin{itemize}
  \item \textbf{Franchissement exact} : PnL Daily à -49\,999 EUR -- pas
        d'alerte. PnL Daily à -50\,001 EUR -- exactement 1 alerte
        \cs{Critical} persistee.
  \item \textbf{Pas de doublon} : deux ticks consécutifs en dessous de la
        limite ne doivent provoquer qu'une seule alerte (pas de duplication).
  \item \textbf{Reset après retour en zone normale} : si le PnL remonte
        au-dessus de la limite puis redescend, une nouvelle alerte doit être
        émise (le système a ``reseté'' son état).
  \item \textbf{Delta limit} : vérifier le déclenchement de l'alerte
        Warning sur le Delta net.
\end{itemize}

\subsubsection{StressTestEngineTests.cs}

\textbf{Cas de test obligatoires :}
\begin{itemize}
  \item \textbf{Parallélisme mesurable} : le temps total de
        \cs{RunAllScenariosAsync} doit être inférieur à la somme
        des temps individuels (marge de 20\%).
  \item \textbf{CombinedShock plus sevère} : le PnL impact du
        CombinedShock doit être inférieur (plus négatif) à
        celui du SpotShock seul pour une position longue call.
  \item \textbf{IProgress reporté} : vérifier que le nombre de
        notifications \cs{IProgress} reçues est égal au nombre de
        scénarios exécutés.
  \item \textbf{Annulation propre} : annuler le \cs{CancellationToken} en
        cours d'exécution ne doit pas lever d'exception non gérée.
\end{itemize}

% ------------------------------------------------------------
\subsection{MarketData/ -- Tests du client API}

\subsubsection{MarketDataServiceTests.cs}

\textbf{Setup} : mocker \cs{HttpMessageHandler} pour simuler les réponses
Alpha Vantage sans appel réseau réel.

\textbf{Cas de test obligatoires :}
\begin{itemize}
  \item \textbf{Cache hit} : deux appels successifs pour le même ticker dans
        la fenêtre d'expiry ne doivent générer qu'un seul appel HTTP.
        Vérifier via \cs{Mock<HttpMessageHandler>} que \cs{SendAsync}
        n'est appelée qu'une fois.
  \item \textbf{Cache miss après expiry} : après expiration du cache
        (manipuler le temps via une interface \cs{ITimeProvider}), un second
        appel doit déclencher un nouvel appel HTTP.
  \item \textbf{Rate limit} : simuler 6 appels simultanés et vérifier que
        le 6ème attend que le token bucket se recharge.
  \item \textbf{HttpRequestException} : simuler une erreur réseau et vérifier
        que l'exception remonte correctement (pas catchée silencieusement).
  \item \textbf{Désérialisation} : vérifier que les string numériques
        Alpha Vantage (\texttt{"182.4500"}) sont correctement converti en
        \cs{decimal}.
\end{itemize}

% ------------------------------------------------------------
\subsection{Contraintes générales du projet Tests}

\begin{tcolorbox}[constraint]
\textbf{FluentAssertions obligatoire :} toutes les assertions utilisent
FluentAssertions. Exemples attendus :
\cs{result.Should().BeApproximately(expected, 0.01)},
\cs{alerts.Should().HaveCount(1)},
\cs{action.Should().NotThrow()},
\cs{result.Should().BeNull()}. \\
Interdiction de \cs{Assert.Equal}, \cs{Assert.True} bruts. \\[4pt]
\textbf{Moq pour toutes les dépendances externes :} aucun accès réseau
réel dans les tests unitaires. Toute dépendance vers DB, API ou système
de fichiers est mockée. Les tests d'intégration (SQLite in-memory)
sont dans des classes séparées suffixees \cs{IntegrationTests}. \\[4pt]
\textbf{Tests async corrects :} toutes les méthodes de test qui appellent du
code async sont \cs{async Task} (jamais \cs{async void}). Toujours
\cs{await} dans le corps du test, jamais \cs{.Result}. \\[4pt]
\textbf{Isolation des tests :} chaque test est indépendant. Pas d'état
partagé entre tests. La base in-memory utilise un nom unique par test
(\cs{Guid.NewGuid()}). Les mocks sont recréés dans chaque méthode
ou dans le constructeur de la classe de test (pas de static).
\end{tcolorbox}

\newpage
% ------------------------------------------------------------
\section{Références entre projets -- Vue complète}
% ------------------------------------------------------------

\begin{center}
\begin{tabularx}{\textwidth}{l l X}
\toprule
\textbf{Projet} & \textbf{Dépend de} & \textbf{Packages NuGet} \\
\midrule
.Core & -- & Aucun \\
.Pricing & .Core & Aucun \\
.Data & .Core & EF Core 8, EF Core SQLite \\
.MarketData & .Core, .Pricing & Polly, System.Text.Json \\
.Risk & .Core, .Pricing & Aucun \\
.Reports & .Core, .Risk & EPPlus, QuestPDF \\
.Dashboard & Tous sauf .Tests & Microsoft.Extensions.DependencyInjection, Hosting \\
.Tests & Tous sauf .Dashboard & xUnit, Moq, FluentAssertions, EF Core InMemory \\
\bottomrule
\end{tabularx}
\end{center}

\begin{tcolorbox}[note]
La règle fondamentale de dépendance : \textbf{jamais de référence
circulaire}. \texttt{Core} ne connaît aucun autre projet. \texttt{Dashboard}
est le seul à avoir une vue globale. Tous les modules communiquent entre eux
via les interfaces définies dans \texttt{Core}. C'est la Clean Architecture
appliquée à un système financier.
\end{tcolorbox}

\end{document}
